% !TEX program = xelatex
% 主文件；第一、二、三章习题分别通过 input 加载。
\documentclass[UTF8,a4paper,11pt,fontset=windows]{ctexart}
\usepackage[margin=2.2cm]{geometry}
\usepackage{amsmath,amssymb,amsthm,enumitem,booktabs,longtable,array}
\usepackage[unicode,hidelinks]{hyperref}
\usepackage{xurl,fancyhdr}
\setlength{\headheight}{15pt}
\pagestyle{fancy}\fancyhf{}
\fancyhead[L]{抽象代数前三章习题详解}\fancyhead[R]{题目、思路与证明}\fancyfoot[C]{\thepage}
\setlength{\parindent}{2em}\setlength{\parskip}{0.35em}
\setlist[enumerate]{itemsep=0.35em,topsep=0.4em}
\renewcommand{\proofname}{证明}
\newcommand{\Z}{\mathbb Z}\newcommand{\Q}{\mathbb Q}
\newcommand{\R}{\mathbb R}\newcommand{\C}{\mathbb C}
\newcommand{\N}{\mathbb N}\newcommand{\F}{\mathbb F}
\DeclareMathOperator{\GL}{GL}\DeclareMathOperator{\SL}{SL}
\DeclareMathOperator{\Aut}{Aut}\DeclareMathOperator{\Inn}{Inn}
\DeclareMathOperator{\lcm}{lcm}\DeclareMathOperator{\ord}{ord}
\DeclareMathOperator{\sgn}{sgn}\DeclareMathOperator{\Ker}{Ker}
\newcommand{\noteref}[1]{\hyperref[note:#1]{\csname notename@#1\endcsname}}
\newcommand{\proofref}[2]{\hyperref[note:#1]{#2}}
\newcommand{\exref}[2]{\hyperref[ex:#1:#2]{第#1章习题#2}}
\newcommand{\registerref}[2]{\expandafter\gdef\csname notename@#1\endcsname{#2}}

\registerref{group}{原笔记第1章第2节，定义 1.2.1（群的定义）}

\registerref{glaws}{原笔记第1章第2节，命题 1.2.3（单位元、逆元与反序取逆）}

\registerref{cancel}{原笔记第1章第2节，命题 1.2.4（方程、消去律）}

\registerref{powers}{原笔记第1章第2节，命题 1.2.5（整数指数法则）}

\registerref{abelianpowers}{原笔记第1章第2节，命题 1.2.7（交换元素的乘积幂）}

\registerref{dihedral}{原笔记第1章第3节，定义 1.3.2（二面体群及其阶数与乘法关系）}

\registerref{cycles}{原笔记第1章第3节，定义 1.3.7（轮换与对换）}

\registerref{cycledecomp}{原笔记第1章第3节，定理 1.3.9（不相交轮换分解）}

\registerref{subgroup}{原笔记第1章第4节，定理 1.4.3（子群判别法）}

\registerref{cosets}{原笔记第1章第4节，命题 1.4.8（陪集相等的判据与陪集划分）}

\registerref{lagrange}{原笔记第1章第4节，定理 1.4.10（拉格朗日定理）}

\registerref{order}{原笔记第1章第4节，命题 1.4.12（元素的阶与幂的判据）}

\registerref{iso}{原笔记第1章第5节，定义 1.5.1（群同构及同构不变量）}

\registerref{cayley}{原笔记第1章第5节，定理 1.5.3（凯莱定理及左正则表示）}

\registerref{sign}{原笔记第1章第4节，命题 1.4.6（符号同态与交错群的阶）}

\registerref{hom}{原笔记第1章第6节，定义 1.6.1（群同态、像与核）}

\registerref{normal}{原笔记第1章第6节，命题 1.6.5（正规子群的等价判据）}

\registerref{setproduct}{原笔记第1章第7节，命题 1.7.2（集合乘积及取逆公式）}

\registerref{quotient}{原笔记第1章第7节，定义 1.7.4（商群及自然同态）}

\registerref{firstiso}{原笔记第1章第7节，定理 1.7.7（群同态基本定理）}

\registerref{ring}{原笔记第1章第8节，定义 1.8.1（环的定义）}

\registerref{ringlaws}{原笔记第1章第8节，命题 1.8.2（环中零元、负元和分配律的运算）}

\registerref{subring}{原笔记第1章第8节，定理 1.8.5（子环判别法）}

\registerref{ringiso}{原笔记第1章第8节，定义 1.8.7（环同构及其性质）}

\registerref{unit}{原笔记第1章第9节，定义 1.9.1（幺环、单位及双侧逆元）}

\registerref{zerodiv}{原笔记第1章第9节，定义 1.9.3（左右零因子、无零因子的消去律）}

\registerref{field}{原笔记第1章第9节，定义 1.9.6（域与体）}

\registerref{quaternion}{原笔记第1章第9节，命题 1.9.9（四元数体的矩阵模型）}

\registerref{ideal}{原笔记第1章第10节，定义 1.10.3（理想的定义）}

\registerref{idealunit}{原笔记第1章第10节，命题 1.10.5（含有单位的理想）}

\registerref{correspondence}{原笔记第2章第1节，定理 2.1.1（子群对应定理）}

\registerref{secondiso}{原笔记第2章第1节，定理 2.1.3（群的第二同构定理）}

\registerref{generated}{原笔记第2章第1节（生成子群与有限生成群）}

\registerref{derived}{原笔记第2章第1节，定义 2.1.6（换位子与换位子群）}

\registerref{abelianization}{原笔记第2章第1节，定理 2.1.7（交换同态像与换位子群）}

\registerref{cyclic}{原笔记第2章第2节，定理 2.2.2（循环群及其子群分类）}

\registerref{powerorder}{原笔记第2章第2节，命题 2.2.3（幂的阶的公式）}

\registerref{coprime}{原笔记第2章第2节，引理 2.2.4（互素阶交换元素乘积的阶）}

\registerref{simple}{原笔记第2章第3节，定义 2.3.1（单群的定义）}

\registerref{threecycles}{原笔记第2章第3节，引理 2.3.3（对换与3-轮换的生成结论）}

\registerref{conjugation}{原笔记第2章第3节，命题 2.3.4（轮换的共轭公式）}

\registerref{an}{原笔记第2章第3节，定理 2.3.6（交错群的单性）}

\registerref{derivedseries}{原笔记第2章第4节，定义 2.4.1（导出列）}

\registerref{solvabledef}{原笔记第2章第4节，定义 2.4.2（可解群的定义）}

\registerref{solvable}{原笔记第2章第4节，命题 2.4.8（可解群的子群、像与扩张）}

\registerref{aut}{原笔记第2章第5节，定义 2.5.1（自同构与自同构群）}

\registerref{cyclicaut}{原笔记第2章第5节，命题 2.5.3（循环群的自同构）}

\registerref{inner}{原笔记第2章第5节，定理 2.5.6（内自同构群与中心商群）}

\registerref{center}{原笔记第2章第5节，定义 2.5.5（群的中心）}

\registerref{autcenterless}{原笔记第2章第5节，定理 2.5.8（无中心群的自同构定理）}

\registerref{complete}{原笔记第2章第5节，定义 2.5.9（完全群）}

\registerref{characteristic}{原笔记第2章第5节，定义 2.5.10（特征子群及其性质）}

\registerref{action}{原笔记第2章第6节，定义 2.6.1（群作用及对应的置换表示）}

\registerref{orbit}{原笔记第2章第6节，定理 2.6.11（轨道稳定子定理）}

\registerref{pcenter}{原笔记第2章第6节，推论 2.6.16（有限$p$群的中心非平凡）}

\registerref{centerizer}{原笔记第2章第6节，定义 2.6.17（元素的中心化子）}

\registerref{classeq}{原笔记第2章第6节，定理 2.6.18（类方程）}

\registerref{normalizer}{原笔记第2章第6节，定义 2.6.20（子群的正规化子）}

\registerref{sylow1}{原笔记第2章第7节，定理 2.7.2（西罗第一定理）}

\registerref{sylowdef}{原笔记第2章第7节，定义 2.7.3（西罗子群的定义）}

\registerref{sylow2}{原笔记第2章第7节，定理 2.7.4（西罗第二定理）}

\registerref{sylow3}{原笔记第2章第7节，定理 2.7.7（西罗第三定理）}

\registerref{internal}{原笔记第2章第8节，定理 2.8.3（内直和分解定理）}

\registerref{free}{原笔记第2章第11节，定理 2.11.7（自由群的泛性质与生成关系）}


\registerref{c3correspondence}{原笔记第3章“环的同态定理”节（环的对应定理，现为源码注释）}
\registerref{c3isomorphism}{原笔记第3章“环的同态定理”节（第三同构定理及其证明，现为源码注释）}
\registerref{c3direct}{原笔记第3章“环的直和”节（直和定义、分量理想与内直和，现为源码注释）}
\registerref{c3product}{原笔记第3章“环的直和”节（理想积的定义与分配律，现为源码注释）}
\registerref{c3comaximal}{原笔记第3章“环的直和”节（互素理想的积等于交，现为源码注释）}
\registerref{c3crt}{原笔记第3章“环的直和”节（中国剩余定理及证明，现为源码注释）}
\registerref{c3principal}{原笔记第3章“环的直和”节（生成理想与主理想的定义及表达式，现为源码注释）}
\registerref{c3anti}{原笔记第3章“环的反同构”节（反同构、反同态与反自同构定义，现为源码注释）}
\registerref{c3conjugate}{原笔记第3章“环的反同构”节（四元数共轭与乘法表，现为源码注释）}
\registerref{c3prime}{原笔记第3章“素理想和极大理想”节（素理想、极大理想及商环判据，现为源码注释）}
\registerref{c3maxexists}{原笔记第3章“素理想和极大理想”节（极大理想存在性与佐恩引理构造，现为源码注释）}
\registerref{c3nil}{原笔记第3章“素理想和极大理想”节（幂零元等于全部素理想之交，现为源码注释）}
\registerref{c3fraction}{原笔记第3章“商域和分式环”节（商域的存在与分式表示，当前正文）}
\registerref{c3poly}{原笔记第3章“交换环上的多项式环”节（多项式环的定义、系数唯一性和代入，当前正文）}
\registerref{c3multi}{原笔记第3章“交换环上的多项式环”节（多元多项式环的定义，当前正文）}
\registerref{c3algebraic}{原笔记第3章“交换环上的多项式环”节（代数关系与代数无关的定义，当前正文）}
\registerref{c3polyquot}{原笔记第3章“交换环上的多项式环”节（系数取商诱导的多项式商环同构，当前正文）}
\registerref{c3degree}{原笔记第3章“整环上的一元多项式环”节（次数法则、整环性与单位，现为源码注释）}
\registerref{c3division}{原笔记第3章“整环上的一元多项式环”节（带余除法与因式定理，现为源码注释）}
\registerref{c3roots}{原笔记第3章“整环上的一元多项式环”节（非零多项式的根数上界，现为源码注释）}
\registerref{c3qroots}{原笔记第3章“整环上的一元多项式环”节（四元数二次方程的共轭根构造，现为源码注释）}
\registerref{c3finite}{原笔记第3章“整环上的一元多项式环”节（有限乘法子群循环、有限域非零元的幂，现为源码注释）}
\registerref{c3minimal}{原笔记第3章“整环上的一元多项式环”节（极小多项式定义与商环刻画，现为源码注释）}
\registerref{c3irred}{原笔记第3章“整环上的一元多项式环”节（不可约定义与极大理想、商域的等价性，现为源码注释）}
\registerref{c3nilideal}{原笔记第3章“整环上的一元多项式环”节（幂零理想定义，现为源码注释）}
\registerref{c3functions1}{原笔记第3章“多项式函数”节（一元多项式函数与形式多项式的对应，现为源码注释）}
\registerref{c3functionsN}{原笔记第3章“多项式函数”节（无限域的多元零函数判据及有限域低次数代表的唯一性，现为源码注释）}

\title{抽象代数前三章习题详解\\\large 第一章1--45题\quad 第二章1--57题\quad 第三章1--68题}
\author{}\date{}
\begin{document}\maketitle
本册收录现有代数笔记抽象代数第一章45题、第二章57题、第三章68题，共170题，保留原题号。原来合写的第一章24--26题分别列出。题面以现有习题文件为基础，必要的限定和已知订正放在相应题目内；证明沿用原笔记的构造、顺序和方法，在原有路线中补齐理由。前两章沿用已经整理的102题；第三章以现有68题及期中复习册已展开的4题为依据补齐证明。

各题按题目、考点与思路、所用结论、证明、复习提示组织。普通证明标题统一为“证明”。题中引用的原笔记定义、定理、命题和证明，以原笔记第一、二章各自的章号、节号和现有统一计数器编号定位；点击引用可跳至文末原笔记定位表，表中列出源文件和源码行号。原笔记正文不在本册重排，也不另建K编号的定理目录。

这里的第一、二、三章均指抽象代数部分，与线性代数部分的章号无关。编号以本次读取的原章节源码顺序为准，不采用当前只启用第三章的主PDF中的临时章号。源码行号仅供本次版本定位，后续手动增删会使其移动。

\paragraph{记号与约定}
群的单位元记为$e$，元素阶记为$\ord(a)$；$H\leq G$表示子群，$N\trianglelefteq G$表示正规子群，$[G:H]$表示左陪集个数。置换乘积从右至左作用。$D_n$为$2n$阶二面体群。$C_m$表示$m$阶循环群，$V_4$表示克莱因四元群。有关定义的正式位置分别见\noteref{group}、\noteref{cosets}、\noteref{order}、\noteref{normal}、\noteref{dihedral}、\noteref{cyclic}。

环不预设有乘法单位元，一般环同态不预设保幺；子环也不要求与母环共用单位元。域指交换的除环，体指不要求交换的除环。零因子必须非零：若$ab=0$且$a,b\ne0$，称$a$为左零因子、$b$为右零因子，沿用\noteref{zerodiv}的左右约定。$E$为单位矩阵，$E_{ij}$为仅第$(i,j)$个位置为1的矩阵单位。

\paragraph{补充工具的归属}
整数带余除法、贝祖等式、有限集计数、数学归纳法，以及线性代数与闭区间紧致性，是原解已经使用的前置工具，不冒充抽象代数正文中另有编号的定理。具体使用时交代条件与计算；不能从原笔记中找到的独立陈述，不赋予原笔记编号。第一章44题逐列计数和第二章47题的有限基构造均在题内展开。第二章12题中沿原笔记补出的附加点引理也在该题内证明。

第三章引用按具名小节和源码行号定位：其中原笔记已注释的小节明确标为“现为源码注释”，不为其编造当前PDF定理编号。已有前两章的引用及题目内容沿用原文件。

\tableofcontents\clearpage

\input{chapter1}

\input{chapter2}

% !TEX root = 抽象代数前三章全部习题详解.tex


\clearpage\section{第3章：环}


本章收录原习题1--68题。涉及$1$、可逆元、素理想和极大理想时，以幺环为背景；交换性及非零条件在需要处说明。幂零元也称零幂元，$r(R)$为幂零根，$J(R)$为全部极大理想之交。第26--29题的若尔当同态统一采用第26题明确给出的三重乘积条件。新补的推导沿现有习题解答展开，不声称已经另行核对原教材书页。


\subsection{习题1：由\texorpdfstring{$1-ab$}{1-ab}可逆推出\texorpdfstring{$1-ba$}{1-ba}可逆}

\label{ex:3:1}

\paragraph{题目}

设环$R$中$1-ab$可逆，证明$1-ba$可逆。

\paragraph{考点与思路}

在非交换环中，不能将 $ab$ 与 $ba$ 直接互换。先记 $c=(1-ab)^{-1}$，再尝试把原来的逆元嵌入 $bca$，构造 $1-ba$ 的逆元。必须分别验证左右两个乘积。

\paragraph{所用结论与原文依据}

\noteref{ringlaws}；\noteref{unit}。原解位置：\path{chapters/Part 2/chapter3-exercises.tex}，第6行起。本题沿用期中复习册第三章已经展开的证明。

% 原题和原解来源同上；新增中间推导为本册补写。

\begin{proof}

由 $c=(1-ab)^{-1}$ 的定义，
\[
(1-ab)c=1,\qquad c(1-ab)=1.
\]
把两式分别展开，得到
\[
c-abc=1,\qquad c-cab=1.
\]
令 $d=1+bca$。先计算 $(1-ba)d$，使用结合律把含 $b$ 和 $a$ 的项提取出来：
\begin{align*}
(1-ba)(1+bca)
  &=1+bca-ba-babca\\
  &=1+b(c-1-abc)a\\
  &=1.
\end{align*}
最后一步用了 $c-abc=1$。

另一方向需要使用另一条逆元等式：
\begin{align*}
(1+bca)(1-ba)
  &=1-ba+bca-bcaba\\
  &=1+b(c-1-cab)a\\
  &=1,
\end{align*}
其中最后一步用了 $c-cab=1$。因此 $d$ 是 $1-ba$ 的双侧逆元，且
\[
\boxed{(1-ba)^{-1}=1+b(1-ab)^{-1}a.}
\]

\end{proof}

\paragraph{复习提示}

记住逆元的构造次序是 $1+bca$。证明中没有交换任意两个因子，因此适用于非交换幺环。将 $a,b$ 对换也得到逆命题，故 $1-ab$ 可逆当且仅当 $1-ba$ 可逆。



\subsection{习题2：右逆不唯一、零因子与非单位的等价性}

\label{ex:3:2}

\paragraph{题目}

设幺环 $R$ 中的元素 $u$ 有右逆。证明下列三个条件等价：
\begin{enumerate}
  \item $u$ 有多于一个右逆；
  \item $u$ 是零因子；
  \item $u$ 不是单位。
\end{enumerate}

\paragraph{考点与思路}

先固定一个右逆 $v$，即 $uv=1$。两个右逆相减后落入方程 $ut=0$ 的解集中；反过来，给 $v$ 加上这样的 $t$ 就能产生新的右逆。条件 $uv=1$ 还排除了非零元素在 $u$ 左侧乘出零的可能。

\paragraph{所用结论与原文依据}

\noteref{ringlaws}；\noteref{unit}。原解位置：\path{chapters/Part 2/chapter3-exercises.tex}，第18行起。本题沿用期中复习册第三章已经展开的证明。

% 原题和原解来源同上；新增中间推导为本册补写。

\begin{proof}

固定 $v\in R$ 满足 $uv=1$。若 $R$ 是只有一个元素的零环，则 $u$ 是单位，只有一个右逆且不是零因子，三个条件均不成立，等价性仍成立。以下设 $1\ne0$，于是 $u\ne0$。

\begin{enumerate}\item 证明条件(1)与条件(2)等价。
若另有 $w\ne v$ 满足 $uw=1$，则 $w-v\ne0$，且
\[
u(w-v)=uw-uv=1-1=0.
\]
因此 $u$ 是左零因子，特别是零因子。

反过来，先注意到：对于任意 $x\in R$，若 $xu=0$，则
\[
x=x\cdot1=x(uv)=(xu)v=0.
\]
所以不存在非零 $x$ 使 $xu=0$。若 $u$ 是零因子，必有非零 $t$ 满足 $ut=0$。于是
\[
u(v+t)=uv+ut=1,
\]
而 $v+t\ne v$，故 $u$ 至少有两个右逆。

\item 证明条件(1)与条件(3)等价。
若 $u$ 是单位，任取右逆 $w$，由 $uw=1$ 得
\[
w=u^{-1}(uw)=u^{-1}.
\]
因而右逆唯一。所以右逆不唯一必推出 $u$ 不是单位。

另一方面，设 $u$ 的右逆唯一。注意
\[
u(vu-1)=(uv)u-u=u-u=0.
\]
因此 $v+(vu-1)$ 也是右逆，因为
\[
u\bigl(v+vu-1\bigr)=uv+u(vu-1)=1.
\]
由右逆的唯一性，$v+vu-1=v$，在加法群中消去 $v$，得到 $vu=1$。结合 $uv=1$，可知 $u$ 是单位。

取这条结论的逆否命题，$u$ 不是单位便推出其右逆不唯一。至此三个条件两两等价。
\end{enumerate}

\end{proof}

\paragraph{复习提示}

这里的前提“已经有一个右逆”不可遗漏。一般幺环中“不是单位”未必是零因子，例如整数环中的 $2$。本题还证明：有右逆的非单位只能以 $ut=0$ 的形式成为零因子，不能有非零 $x$ 满足 $xu=0$。



\subsection{习题3：右逆不唯一便有无穷多个右逆}

\label{ex:3:3}

\paragraph{题目}

若$u$有多于一个右逆，证明它有无穷多个右逆。

\paragraph{考点与思路}

由 $uv=1$ 但 $vu\ne1$ 构造非零元素 $e=1-vu$。它满足 $ue=0$、$ev=0$，适合用来给已有的右逆 $v$ 添加不同的修正项。不能仅写 $v+ne$，因为正特征下这些元素可能重复。

\paragraph{所用结论与原文依据}

\noteref{ringlaws}；\noteref{unit}。原解位置：\path{chapters/Part 2/chapter3-exercises.tex}，第31行起。本题沿用期中复习册第三章已经展开的证明。

% 原题和原解来源同上；新增中间推导为本册补写。

\begin{proof}

取一个右逆 $v$，使 $uv=1$。由\exref{3}{2}，$u$ 不是单位。若 $vu=1$，则 $v$ 同时是左逆和右逆，$u$ 便是单位，矛盾。因此
\[
e=1-vu\ne0.
\]
由 $uv=1$，有
\begin{align*}
ue&=u(1-vu)=u-(uv)u=0,\\
ev&=(1-vu)v=v-v(uv)=0.
\end{align*}

对每个整数 $k\ge0$，令
\[
x_k=v+eu^k,
\]
这里 $u^0=1$。因为 $ue=0$，
\[
ux_k=uv+(ue)u^k=1,
\]
所以每个 $x_k$ 都是 $u$ 的右逆。

下面证明它们两两不同。先由 $uv=1$ 归纳得到，对每个整数 $m\ge0$，
\[
u^mv^m=1.
\]
归纳步骤为
\[
u^{m+1}v^{m+1}=u^m(uv)v^m=u^mv^m=1.
\]
于是当 $0\le i<j$ 时，
\[
u^iv^j=(u^iv^i)v^{j-i}=v^{j-i}.
\]

假设 $x_i=x_j$，其中 $i<j$。消去 $v$ 后，$eu^i=eu^j$。两边右乘 $v^j$，得到
\[
ev^{j-i}=e.
\]
但 $j-i\ge1$ 且 $ev=0$，所以
\[
ev^{j-i}=(ev)v^{j-i-1}=0,
\]
从而 $e=0$，与构造矛盾。因此 $x_0,x_1,x_2,\ldots$ 两两不同，$u$ 有无穷多个右逆。

\end{proof}

\paragraph{复习提示}

证明的核心是“构造右逆”和“证明构造出的元素互异”两个步骤。等式 $u^mv^m=1$ 来自逐对消去中间的 $uv$，绝不意味着 $u,v$ 可以交换。本题也说明有限幺环中，有右逆的元素必为单位。



\subsection{习题4：幂零元全体是否为理想}

\label{ex:3:4}

\paragraph{题目}

证明交换环中全部幂零元构成理想，并判断非交换环中的情形。

\paragraph{考点与思路}

交换情形按原解展开二项式，并分别检查减法和吸收性；非交换情形保留两个矩阵单位的反例。

\paragraph{所用结论与原文依据}

\noteref{ringlaws}；\noteref{ideal}；\noteref{c3nil}。原解位置：\path{chapters/Part 2/chapter3-exercises.tex}，第42行起。本题在现有习题解答的思路内补齐推导。

% 原题和原解来源同上；新增中间推导为本册补写。

\begin{proof}

令$N=\{a\in R:a^m=0\text{ 对某个 }m\ge1\}$。因为$0\in N$，它非空。
任取$a,b\in N$，选$m,n\ge1$使$a^m=b^n=0$。因$R$交换，二项式展开的每一项形如$a^i(-b)^j$，其中$i+j=m+n-1$。若$i<m$且$j<n$，则$i+j\le m+n-2$，矛盾。因此每一项均含$a^m$或$b^n$，从而
\[
(a-b)^{m+n-1}=0.
\]
这说明$a-b\in N$。若环无幺元，两个端项分别直接理解为$a^{m+n-1}$和$(-b)^{m+n-1}$，无须使用零次幂。

任取$r\in R$，交换性给出$(ra)^m=r^ma^m=0$，故$ra=ar\in N$。所以$N$是理想。

在非交换环$M_2(\Z)$中，取$A=E_{12}$、$B=E_{21}$。有
\[
A^2=B^2=0,\qquad (A+B)^2=E.
\]
于是$(A+B)^{2k}=E\ne0$，$(A+B)^{2k+1}=A+B\ne0$，所以$A+B$不幂零。幂零元全体甚至不对加法封闭，因而一般不是理想。

\end{proof}

\paragraph{复习提示}

二项式展开中将同类乘积合并使用了交换性；不能把该证明直接移用于非交换环。



\subsection{习题5：素特征的幂公式与一般环的特征}

\label{ex:3:5}

\paragraph{题目}

设$F$的特征为素数$p$，证明$(a+b)^p=a^p+b^p$。并将特征的讨论推广到幺环及不必有幺元的环。

\paragraph{考点与思路}

先证明中间二项式系数被素数整除，再把特征解释为加法群的共同湮灭整数。

\paragraph{所用结论与原文依据}

\noteref{ring}；\noteref{ringlaws}；\noteref{field}。原解位置：\path{chapters/Part 2/chapter3-exercises.tex}，第51行起。本题在现有习题解答的思路内补齐推导。

% 原题和原解来源同上；新增中间推导为本册补写。

\begin{proof}

\begin{enumerate}
\item 设$\operatorname{char}F=p$。对$1\le k\le p-1$，整数等式
\[
k\binom pk=p\binom{p-1}{k-1}
\]
及$p\nmid k$说明$p\mid\binom pk$。因此这些整数系数作用于$F$中的元素时均为零。$F$的乘法交换，故
\[
(a+b)^p=a^p+\sum_{k=1}^{p-1}\binom pk a^kb^{p-k}+b^p=a^p+b^p.
\]
\item 对任意幺环$R$，若存在正整数$n$使$n1=0$，就以最小的这样的$n$为特征；不存在时特征为$0$。这里$n1$表示加法意义下的$n$倍。由$(n1)r=nr$，条件$n1=0$等价于$nR=0$。

一般幺环的正特征未必为素数，例如$\Z/(6)$的特征为$6$。若非零幺环无零因子，设其正特征$n=st$，其中$1<s,t<n$，则
\[
(s1)(t1)=n1=0,
\]
而最小性说明$s1,t1$均非零，矛盾。因此此时$n$为素数。
\item 对不要求有幺元的环，定义特征为使$nr=0$对所有$r\in R$成立的最小正整数$n$；若无此$n$则为$0$。这一定义只用加法，不需要元素$1$。不能一般改成“某个非零元素的阶”。

若环交换且特征为素数$p$，第一项的二项式证明仍然适用；端项写为$a^p,b^p$即可。非交换环即使特征为$p$，也不能由此得到该公式，因为混合乘积不能按二项式系数合并。
\end{enumerate}

\end{proof}

\paragraph{复习提示}

一般幺环的特征可以是合数；素数结论需要非零且无零因子的条件。



\subsection{习题6：矩阵环的全部双边理想}

\label{ex:3:6}

\paragraph{题目}

设$R$为幺环，$n\ge1$，$I\lhd R$。证明$M_n(I)\lhd M_n(R)$，且$M_n(R)$的每个双边理想均唯一地如此得到。原题未明写的幺元条件在此列出；无幺情形见证明末的反例。

\paragraph{考点与思路}

用矩阵单位取出一个分量并把它搬到任意位置，再重组整个矩阵。

\paragraph{所用结论与原文依据}

\noteref{ideal}；\noteref{unit}。原解位置：\path{chapters/Part 2/chapter3-exercises.tex}，第62行起。本题在现有习题解答的思路内补齐推导。

% 原题和原解来源同上；新增中间推导为本册补写。

\begin{proof}

这里$R$为幺环，$n\ge1$，理想均为双边理想。若$I\lhd R$，则$M_n(I)$对减法封闭；对$A\in M_n(R)$、$B\in M_n(I)$，乘积$AB,BA$的每个分量都是含有$I$中一个因子的有限和，仍属于$I$。故$M_n(I)\lhd M_n(R)$。

反过来，设$J\lhd M_n(R)$。令
\[
I=\{a\in R:aE_{11}\in J\}.
\]
零矩阵属于$J$，故$0\in I$。若$a,b\in I$、$r\in R$，则
\begin{align*}
(a-b)E_{11}&=aE_{11}-bE_{11}\in J,\\
(ra)E_{11}&=(rE_{11})(aE_{11})\in J,\\
(ar)E_{11}&=(aE_{11})(rE_{11})\in J.
\end{align*}
因此$I\lhd R$。

对$A=(a_{ij})\in J$，矩阵单位给出
\[
E_{1i}AE_{j1}=a_{ij}E_{11}\in J,
\]
所以每个$a_{ij}\in I$，即$J\subseteq M_n(I)$。反之，对$a\in I$及任意$i,j$，
\[
E_{i1}(aE_{11})E_{1j}=aE_{ij}\in J.
\]
将这些矩阵相加，$M_n(I)$中的每个矩阵都属于$J$，故$J=M_n(I)$。由$a\in I\Longleftrightarrow aE_{11}\in J$，$I$也唯一。

幺元假设保证所用矩阵单位属于$M_n(R)$。无幺环中不能直接断言此结论：取加法群为$\Z$、乘法恒为零的环$R$，则$M_2(R)$也是零乘法环，任一加法子群均为理想；子群$\Z E_{11}$并非任何$M_2(I)$。

\end{proof}

\paragraph{复习提示}

要证明全部理想的分类，既要从任意理想恢复$I$，也要验证任意合适的$I$确实给出理想。



\subsection{习题7：理想积的分配律}

\label{ex:3:7}

\paragraph{题目}

若$I,H,N$为$R$的理想，证明$(H+N)I=HI+NI$及$I(H+N)=IH+IN$。

\paragraph{考点与思路}

按理想积的有限和定义，逐项展开，再分别证明反向包含。

\paragraph{所用结论与原文依据}

\noteref{c3product}；\noteref{ringlaws}。原解位置：\path{chapters/Part 2/chapter3-exercises.tex}，第75行起。本题在现有习题解答的思路内补齐推导。

% 原题和原解来源同上；新增中间推导为本册补写。

\begin{proof}

按定义，$HI$是全部有限和$\sum h_ji_j$组成的集合，其中$h_j\in H,i_j\in I$；其余理想积也按同样方式理解。

任取$x\in(H+N)I$，写作$x=\sum_j(h_j+n_j)i_j$。分配律给出
\[
x=\sum_jh_ji_j+\sum_jn_ji_j\in HI+NI.
\]
反过来，因$H,N\subseteq H+N$，每个$h_ji_j$和$n_ki_k$都是$(H+N)I$的生成乘积；该集合对加法封闭，故$HI+NI\subseteq(H+N)I$。于是$(H+N)I=HI+NI$。

对另一等式，任取$y=\sum_ji_j(h_j+n_j)\in I(H+N)$，则
\[
y=\sum_ji_jh_j+\sum_ji_jn_j\in IH+IN.
\]
而$H,N\subseteq H+N$分别给出$IH,IN\subseteq I(H+N)$，故$IH+IN\subseteq I(H+N)$。得到$I(H+N)=IH+IN$。两式均不要求乘法交换。

\end{proof}

\paragraph{复习提示}

理想积的元素是有限个乘积之和，不一定是单个乘积。



\subsection{习题8：满同态下的元素性质与整环性质}

\label{ex:3:8}

\paragraph{题目}

设 $\eta:R\twoheadrightarrow R'$ 为保幺满同态。判断下列命题是否成立，并证明或举反例：
\begin{enumerate}
  \item[(i)] 幂零元、幂等元的像仍是相应元素；
  \item[(ii)] 零因子的像仍是零因子；
  \item[(iii)] 若 $R$ 为整环，则 $R'$ 为整环；
  \item[(iv)] 若 $R'$ 为整环，则 $R$ 为整环；
  \item[(v)] 单位的像仍是单位；
  \item[(vi)] 像为单位的元素必为单位。
\end{enumerate}

\paragraph{考点与思路}

同态保留由加法、乘法和 $0,1$ 写出的等式，却未必保留“非零”“不存在非零零乘积”等条件。零因子的定义涉及非零元素，所以不能只把 $ab=0$ 两边取像就认定像仍为零因子。

\paragraph{所用结论与原文依据}

\noteref{ringiso}；\noteref{ideal}。原解位置：\path{chapters/Part 2/chapter3-exercises.tex}，第82行起。本题沿用期中复习册第三章已经展开的证明。

% 原题和原解来源同上；新增中间推导为本册补写。

\begin{proof}

\begin{enumerate}[label=(\roman*)]\item 成立。
若 $a$ 是幂零元，取正整数 $m$ 使 $a^m=0$。由同态保乘法，
\[
\eta(a)^m=\eta(a^m)=\eta(0)=0'.
\]
故 $\eta(a)$ 仍为幂零元。若 $e^2=e$，则
\[
\eta(e)^2=\eta(e^2)=\eta(e),
\]
故 $\eta(e)$ 仍为幂等元。两项结论都不要求满射或保幺。

\item 不成立。
考虑映射
\[
\eta:\Z/6\Z\longrightarrow\Z/3\Z,\qquad
\overline{k}_{6}\longmapsto\overline{k}_{3}.
\]
若 $k\equiv\ell\pmod6$，则 $3\mid k-\ell$，所以像不依赖代表元，映射良定义。模运算保留加法和乘法，且 $\eta(\overline1_6)=\overline1_3$；每个 $\overline{k}_3$ 都是 $\overline{k}_6$ 的像，故这是保幺满同态。

在源环中，$\overline4_6\ne0$、$\overline3_6\ne0$，且
\[
\overline4_6\,\overline3_6=\overline{12}_6=\overline0_6,
\]
所以 $\overline4_6$ 是零因子。它的像却是 $\overline1_3$。若 $\overline1_3 y=0$，则 $y=0$，故该像不是零因子。这一反例甚至将零因子映成了非零单位。

\item 不成立。
取自然映射
\[
\pi:\Z\twoheadrightarrow\Z/6\Z,\qquad k\longmapsto\overline{k}_6.
\]
该映射保加法、乘法和单位元，并且满射。源环 $\Z$ 是整环，但目标中
\[
\overline2_6\ne0,\qquad \overline3_6\ne0,\qquad
\overline2_6\,\overline3_6=0,
\]
所以 $\Z/6\Z$ 不是整环。

\item 不成立。
取映射
\[
\rho:\Z/6\Z\twoheadrightarrow\Z/2\Z,\qquad
\overline{k}_{6}\longmapsto\overline{k}_{2}.
\]
因为 $2\mid6$，同余模 $6$ 蕴含同余模 $2$，映射良定义；模运算说明它保加法、乘法及单位元，每个模 $2$ 的剩余类都有原像，所以满射。

目标 $\Z/2\Z$ 只有两个元素，唯一非零元素是 $\overline1_2$，且 $\overline1_2\,\overline1_2=\overline1_2\ne0$，因此它是整环。源环却有上一项给出的非零零乘积，故不是整环。

\item 成立。
若 $u$ 是单位，令 $v=u^{-1}$，则 $uv=vu=1$。取像得
\begin{align*}
\eta(u)\eta(v)&=\eta(uv)=\eta(1)=1',\\
\eta(v)\eta(u)&=\eta(vu)=\eta(1)=1'.
\end{align*}
因此 $\eta(u)$ 是单位，且
\[
\eta(u)^{-1}=\eta(u^{-1}).
\]
这一结论只需要保幺同态，不需要满射。

\item 不成立。
取自然保幺满同态
\[
\sigma:\Z\twoheadrightarrow\Z/5\Z,\qquad
k\longmapsto\overline{k}_5.
\]
整数 $2$ 不是 $\Z$ 中的单位：若 $2m=1$，则偶数等于奇数，故没有整数解。它的像 $\overline2_5$ 却满足
\[
\overline2_5\,\overline3_5=\overline6_5=\overline1_5,
\]
因而是单位。这说明像可逆不保证原元素可逆。
\end{enumerate}

\end{proof}

\paragraph{复习提示}

六项判断依次为：\textbf{对、错、错、错、对、错}。构造剩余类环之间的映射时，先检验模数之间的整除关系，确保“按原整数取剩余类”的规则良定义。证明零因子的反例时，须同时写明两个因子非零；证明单位的反例时，须给出像的实际逆元。



\subsection{习题9：有理系数四元数构成子体}

\label{ex:3:9}

\paragraph{题目}

证明有理系数四元数$H_0=\{a_0+a_1I+a_2J+a_3K:a_i\in\mathbb Q\}$是$H$的子体。

\paragraph{考点与思路}

沿用四元数乘法表检查封闭性，再用共轭除以平方和求逆。

\paragraph{所用结论与原文依据}

\noteref{quaternion}；\noteref{c3conjugate}；\noteref{field}。原解位置：\path{chapters/Part 2/chapter3-exercises.tex}，第99行起。本题在现有习题解答的思路内补齐推导。

% 原题和原解来源同上；新增中间推导为本册补写。

\begin{proof}

$H_0$含$0,1$，且对减法封闭。将两个有理系数四元数的乘积按$1,I,J,K$展开，由乘法表可知每个新系数都是原系数的整数系数和积，仍为有理数。因此$H_0$是$H$的含幺子环。

任取$0\ne q=a_0+a_1I+a_2J+a_3K\in H_0$，定义$\bar q=a_0-a_1I-a_2J-a_3K$。不同虚基元的交叉项成对抵消，得到
\[
q\bar q=\bar q q=a_0^2+a_1^2+a_2^2+a_3^2=:s.
\]
至少一个$a_i$非零，且这些系数属于$\Q\subset\R$，所以$s>0$。特别地，$s$是非零有理数，故
\[
q^{-1}=s^{-1}\bar q\in H_0.
\]
左右乘$q$均得到$1$，所以每个非零元素都有逆，$H_0$是$H$的子体。这里“体”指除环；$IJ=K$、$JI=-K$说明它不是交换域。

\end{proof}

\paragraph{复习提示}

平方和非零的理由来自系数位于实数中；在任意域上不能直接使用正性。



\subsection{习题10：实四元数体的中心}

\label{ex:3:10}

\paragraph{题目}

确定四元数体$H$的中心。

\paragraph{考点与思路}

先令一般四元数分别与$I,J$交换，以系数唯一性逐次排除虚部。

\paragraph{所用结论与原文依据}

\noteref{quaternion}；\noteref{c3conjugate}。原解位置：\path{chapters/Part 2/chapter3-exercises.tex}，第111行起。本题在现有习题解答的思路内补齐推导。

% 原题和原解来源同上；新增中间推导为本册补写。

\begin{proof}

环的中心定义为$Z(H)=\{q\in H:qr=rq\text{ 对每个 }r\in H\}$。任取$q=a+bI+cJ+dK\in Z(H)$。由乘法表，
\begin{align*}
qI&=-b+aI+dJ-cK,\\
Iq&=-b+aI-dJ+cK.
\end{align*}
由$qI=Iq$及$1,I,J,K$的实线性无关性，有$2d=2c=0$，故$c=d=0$。

此时$q=a+bI$，又
\[
qJ=aJ+bK,\qquad Jq=aJ-bK.
\]
因$qJ=Jq$，得$2b=0$，故$b=0$。所以$Z(H)\subseteq\R1$。
反过来，任一实标量$a1$与四元数的各基元交换，对分配律展开后的任意四元数也交换。因此
\[
Z(H)=\R1\cong\R.
\]

\end{proof}

\paragraph{复习提示}

必要性只需与两个特定元素交换；充分性仍须说明实标量与所有四元数交换。



\subsection{习题11：与四元数基元不交换的元素}

\label{ex:3:11}

\paragraph{题目}

求$H$中与$I$不交换的全部元素。

\paragraph{考点与思路}

计算交换差$I q-q I$，再用$J,K$的线性无关性判定何时非零。

\paragraph{所用结论与原文依据}

\noteref{quaternion}。原解位置：\path{chapters/Part 2/chapter3-exercises.tex}，第118行起。本题在现有习题解答的思路内补齐推导。

% 原题和原解来源同上；新增中间推导为本册补写。

\begin{proof}

写$q=a+bI+cJ+dK$，其中$a,b,c,d\in\R$。乘法表给出
\begin{align*}
Iq&=-b+aI-dJ+cK,\\
qI&=-b+aI+dJ-cK,\\
Iq-qI&=2cK-2dJ.
\end{align*}
由于$J,K$在实数上独立，$Iq-qI=0$当且仅当$c=d=0$。所以全部与$I$不交换的元素恰为
\[
\{a+bI+cJ+dK:a,b,c,d\in\R,\ (c,d)\ne(0,0)\}.
\]

\end{proof}

\paragraph{复习提示}

“不交换”要求交换差非零，而不是要求四个系数都非零。



\subsection{习题12：共轭不变的四元数子体}

\label{ex:3:12}

\paragraph{题目}

设$S$是$H$的子体，且对所有$0\ne d\in H$有$dSd^{-1}\subseteq S$。证明$S=H$或$S$包含于中心。

\paragraph{考点与思路}

保留原解中$x$与$x+1$的两个共轭，比较后迫使它们相等；再以$x+y$处理子体内部的元素。

\paragraph{所用结论与原文依据}

\noteref{field}；\noteref{c3conjugate}；\exref{3}{10}。原解位置：\path{chapters/Part 2/chapter3-exercises.tex}，第129行起。本题在现有习题解答的思路内补齐推导。

% 原题和原解来源同上；新增中间推导为本册补写。

\begin{proof}

若$S=H$则结论成立。以下设$S\ne H$，固定$x\in H\setminus S$。子体$S$含$0,1$，所以$x\ne0$且$x+1\ne0$。任取$s\in S$，题设保证
\[
a=xsx^{-1}\in S,\qquad b=(x+1)s(x+1)^{-1}\in S.
\]
右乘$x$及$x+1$分别得$xs=ax$及$xs+s=bx+b$。相减得
\[
(a-b)x=b-s.
\]
若$a\ne b$，则$a-b$在子体$S$中可逆，从左侧消去得$x=(a-b)^{-1}(b-s)\in S$，矛盾。故$a=b$，上述等式进而给出$b=s$，于是$xs=sx$。

由于$x$在$H\setminus S$中可任取，子体外的每个元素都与$s$交换。对任意$y\in S$，固定一个$x\notin S$；若$x+y\in S$则$x=(x+y)-y\in S$，矛盾。因此$x+y\notin S$，它也与$s$交换。于是
\[
ys=(x+y)s-xs=s(x+y)-sx=sy.
\]
所以$s$与$S$内外的所有元素均交换，即$s\in Z(H)$。$s$任意，故$S\subseteq Z(H)=\R1$。

\end{proof}

\paragraph{复习提示}

最后以$x+y$相减，才能从子体外的交换性推出子体内部的交换性。



\subsection{习题13：互素理想与理想积的包含}

\label{ex:3:13}

\paragraph{题目}

设$R$为幺环，$I,H,N\lhd R$。若$HN\subseteq I$且$I+H=R$，证明$N\subseteq I$。这里补明原解已经使用的幺元条件。

\paragraph{考点与思路}

由$I+H=R$把$1$分成两项，再乘以$N$中的任意元素。

\paragraph{所用结论与原文依据}

\noteref{ideal}；\noteref{c3product}。原解位置：\path{chapters/Part 2/chapter3-exercises.tex}，第142行起。本题在现有习题解答的思路内补齐推导。

% 原题和原解来源同上；新增中间推导为本册补写。

\begin{proof}

这里$R$为幺环。由$I+H=R$，可取$i\in I,h\in H$使$1=i+h$。任取$n\in N$，有
\[
n=1n=(i+h)n=in+hn.
\]
因$I$为理想，$in\in I$；又$h\in H,n\in N$，所以$hn\in HN\subseteq I$。故$n\in I$。$n$任意，得到$N\subseteq I$。

此证明不需要交换性，但需要幺元。若允许无幺环，可在非零的零乘法环中取$I=0,H=N=R$，此时$HN=0\subseteq I$、$I+H=R$，却$N\nsubseteq I$。

\end{proof}

\paragraph{复习提示}

用到$1=i+h$的题目须明确母环有幺元。



\subsection{习题14：互素理想之积包含其交中的理想}

\label{ex:3:14}

\paragraph{题目}

设$I,H,N$为交换幺环$R$的理想，$I\subseteq H,N$且$H+N=R$，证明$HN\supseteq I$。补明原解中使用的幺元条件。

\paragraph{考点与思路}

由$1=h+n$把$I$中的元素拆成两个属于$HN$的项。

\paragraph{所用结论与原文依据}

\noteref{c3comaximal}；\noteref{c3product}。原解位置：\path{chapters/Part 2/chapter3-exercises.tex}，第153行起。本题在现有习题解答的思路内补齐推导。

% 原题和原解来源同上；新增中间推导为本册补写。

\begin{proof}

这里$R$为交换幺环。由$H+N=R$，取$h\in H,n\in N$使$h+n=1$。任取$i\in I$。因$I\subseteq H\cap N$，有$i\in H$且$i\in N$，所以
\[
i=i(h+n)=ih+in=hi+in\in HN.
\]
其中$hi\in HN$使用$h\in H,i\in N$，$in\in HN$使用$i\in H,n\in N$。故$I\subseteq HN$。

这里$\supseteq$允许相等：取$R=F\oplus F$、$I=0$、$H=F\oplus0$、$N=0\oplus F$，则$H+N=R$而$HN=0=I$。
若删去幺元假设，非零零乘法环中取$I=H=N=R$就使结论失败。

\end{proof}

\paragraph{复习提示}

交换性在将$ih$改写为$hi$时使用；包含关系不能无条件加强为严格包含。



\subsection{习题15：只有恒等反自同构的矩阵子环}

\label{ex:3:15}

\paragraph{题目}

在$M_n(F)$（$n>1$）中找一个子环$R$，使$R$除恒等同构外没有其他反自同构。

\paragraph{考点与思路}

沿原解取素子域的标量矩阵环；先证明自同构必须固定幺元，再固定全部素域元素。

\paragraph{所用结论与原文依据}

\noteref{c3anti}；\noteref{unit}；\noteref{field}。原解位置：\path{chapters/Part 2/chapter3-exercises.tex}，第164行起。本题在现有习题解答的思路内补齐推导。

% 原题和原解来源同上；新增中间推导为本册补写。

\begin{proof}

令$F_0$为$F$的素子域，即$\operatorname{char}F=0$时的$\Q$，或特征为$p$时的$\F_p$，并取
\[
R=\{cE_n:c\in F_0\}\subseteq M_n(F).
\]
标量矩阵的加法、乘法对应于$F_0$的运算，故$R$是交换子环，单位元为$E_n$。

设$\sigma:R\to R$为反自同构。因$R$交换，$\sigma(AB)=\sigma(B)\sigma(A)=\sigma(A)\sigma(B)$，故它也是自同构。即使不预设保幺，满射性也迫使$\sigma(E_n)=E_n$：对任意$B=\sigma(A)$，都有$\sigma(E_n)B=B\sigma(E_n)=B$，再利用单位元唯一性。

加法性说明$\sigma$固定每个整数倍$mE_n$。正特征时这些就是全部$F_0$标量。零特征时，对非零整数$m$，由$\sigma$保乘法，$\sigma((mE_n)^{-1})=(mE_n)^{-1}$，继而固定每个有理数标量。因此两种情形都只有$\sigma=\mathrm{id}_R$。

\end{proof}

\paragraph{复习提示}

题目只要求一个子环，素子域标量环已经满足要求；不必让子环等于整个矩阵环。



\subsection{习题16：整数下三角矩阵环的全部理想}

\label{ex:3:16}

\paragraph{题目}

求$R=\left\{\begin{pmatrix}a&0\\b&c\end{pmatrix}:a,b,c\in\mathbb Z\right\}$的全部理想。

\paragraph{考点与思路}

用对角矩阵单位分离三个位置，再用$E_{21}$比较对应的三个整数理想。

\paragraph{所用结论与原文依据}

\noteref{ideal}；\noteref{cyclic}。原解位置：\path{chapters/Part 2/chapter3-exercises.tex}，第171行起。本题在现有习题解答的思路内补齐推导。

% 原题和原解来源同上；新增中间推导为本册补写。

\begin{proof}

设$J\lhd R$。对$A=\left(\begin{smallmatrix}a&0\\b&c\end{smallmatrix}\right)\in J$，有
\[
E_{11}AE_{11}=aE_{11},\quad E_{22}AE_{11}=bE_{21},\quad E_{22}AE_{22}=cE_{22},
\]
这些矩阵都在$J$中。分别令$A_1,B_1,C_1$为满足$aE_{11}\in J,bE_{21}\in J,cE_{22}\in J$的整数集合。它们都是$\Z$的加法子群，故分别唯一地形如$m\Z,k\Z,n\Z$，其中$m,k,n\ge0$。

分离公式和加法封闭性共同说明
\[
J=I_{m,k,n}:=\left\{\begin{pmatrix}a&0\\b&c\end{pmatrix}:a\in m\Z,\ b\in k\Z,\ c\in n\Z\right\}.
\]
又因$E_{21}(aE_{11})=aE_{21}$、$(cE_{22})E_{21}=cE_{21}$，必有
\[
m\Z\subseteq k\Z,\qquad n\Z\subseteq k\Z.
\]

反过来，取任意满足上述两包含关系的$m,k,n$。$I_{m,k,n}$对减法封闭。对$X=\left(\begin{smallmatrix}r&0\\s&t\end{smallmatrix}\right)\in R$和$A\in I_{m,k,n}$，
\[
XA=\begin{pmatrix}ra&0\\sa+tb&tc\end{pmatrix},\qquad
AX=\begin{pmatrix}ar&0\\br+cs&ct\end{pmatrix}.
\]
两个对角分量分别仍在$m\Z,n\Z$中；左下分量中$a,c\in k\Z$，故所有项均在$k\Z$中。因此$XA,AX\in I_{m,k,n}$，它确为双边理想。

所以全部理想恰是上述$I_{m,k,n}$。参数条件也可写为$k\mid m$、$k\mid n$，其中$0\mid m$仅当$m=0$；用集合包含写法可统一涵盖零参数。

\end{proof}

\paragraph{复习提示}

先取出各分量，才能排除不同位置的系数之间还存在额外耦合关系。



\subsection{习题17：剩余类矩阵环的单位群阶数}

\label{ex:3:17}

\paragraph{题目}

确定$M_2(\mathbb Z/(n))$的单位群的阶。

\paragraph{考点与思路}

保留中国剩余分解、模$p$约化、两列计数的路线，补齐可逆性提升的理由。

\paragraph{所用结论与原文依据}

\noteref{c3crt}；\noteref{unit}；\exref{1}{44}。原解位置：\path{chapters/Part 2/chapter3-exercises.tex}，第184行起。本题在现有习题解答的思路内补齐推导。

% 原题和原解来源同上；新增中间推导为本册补写。

\begin{proof}

设$n\ge2$，写$n=\prod_{i=1}^r p_i^{e_i}$。中国剩余定理按矩阵分量给出环同构
\[
M_2(\Z/(n))\cong\bigoplus\limits_{i=1}^r M_2(\Z/(p_i^{e_i})).
\]
分量逐一可逆当且仅当整个元组可逆，所以其单位群是各$\GL_2(\Z/(p_i^{e_i}))$的直积，阶数相乘。

固定$p^e$。交换环上的二阶矩阵$A$可逆当且仅当$\det A$为单位：必要性由$\det A\det A^{-1}=1$得到；充分性由$A\operatorname{adj}A=\operatorname{adj}A\,A=(\det A)E$得到。整数剩余类$a\bmod p^e$为单位当且仅当$p\nmid a$，因为互素时贝祖等式给出逆；被$p$整除时不能与任何整数相乘得到$1$模$p^e$。

因此$A$在模$p^e$下可逆，当且仅当其模$p$像可逆。每个模$p$的矩阵有$p^{4(e-1)}$个提升，四个分量均可独立选择；一个可逆模$p$矩阵的所有这些提升都可逆。于是
\[
|\GL_2(\Z/(p^e))|=p^{4(e-1)}|\GL_2(\F_p)|.
\]
在$\F_p^2$中，可逆矩阵的两列须为有序基。第一列可选任意非零向量，共$p^2-1$种；它张成的直线有$p$个向量，第二列须避开此直线，共$p^2-p$种。因此
\[
\boxed{|\GL_2(\Z/(n))|=\prod_{p^e\parallel n}p^{4(e-1)}(p^2-1)(p^2-p).}
\]
若$n=1$，环只有一个元素，且$0=1$；按单位的双侧逆定义，单位群为单元素群，阶为$1$，也与空乘积约定一致。

\end{proof}

\paragraph{复习提示}

约化后的矩阵可逆，才能保证它的每个提升可逆；理由是行列式为单位的判据。



\subsection{习题18：有限无零因子交换环自动有幺且为域}

\label{ex:3:18}

\paragraph{题目}

设$R\ne0$为交换环，不预设幺元，且无零因子、非零元素只有有限个。证明$R$是有限除环。原题应补明非零条件；事实上结论是有限域。

\paragraph{考点与思路}

不预设幺元，先由有限集上的乘法单射得到一个对所有元素有效的单位元。

\paragraph{所用结论与原文依据}

\noteref{zerodiv}；\noteref{field}。原解位置：\path{chapters/Part 2/chapter3-exercises.tex}，第199行起。本题在现有习题解答的思路内补齐推导。

% 原题和原解来源同上；新增中间推导为本册补写。

\begin{proof}

补明$R\ne0$；否则零环满足原题的有限性和无非零零因子条件，却不是除环。固定$0\ne a\in R$。若$ax=ay$，则$a(x-y)=0$；无零因子性使$x=y$。故$L_a:x\mapsto ax$是有限集合$R$上的单射，因而满射。

取$e\in R$使$ae=a$。任给$b\in R$，由满射性写$b=ac$。交换性和结合律给出
\[
eb=e(ac)=(ea)c=(ae)c=ac=b,\qquad be=eb=b.
\]
因此$e$是整个$R$的双侧单位元。由$a\ne0$、$ae=a$知$e\ne0$。

再取任意$0\ne b\in R$，同样的消去论证说明$L_b$满射，故存在$c\in R$使$bc=e$。交换性给出$cb=e$。所以每个非零元都有双侧逆元，$R$是有限域，特别是有限除环。

\end{proof}

\paragraph{复习提示}

构造$ae=a$后，还须证明$e$对任意$b$都是单位；$R\ne0$也不能遗漏。



\subsection{习题19：理想关系不具传递性}

\label{ex:3:19}

\paragraph{题目}

若$N\lhd R$而$I\lhd N$，问$I$是否必为$R$的理想？

\paragraph{考点与思路}

沿用严格上三角矩阵的反例，写明每个乘积落在哪里及哪一个乘积破坏吸收性。

\paragraph{所用结论与原文依据}

\noteref{ideal}。原解位置：\path{chapters/Part 2/chapter3-exercises.tex}，第206行起。本题在现有习题解答的思路内补齐推导。

% 原题和原解来源同上；新增中间推导为本册补写。

\begin{proof}

设$F$为域，取
\[
R=\left\{\begin{pmatrix}0&a&b\\0&0&c\\0&0&0\end{pmatrix}:a,b,c\in F\right\},\quad
N=FE_{12}+FE_{13},\quad I=FE_{12}.
\]
对$X=aE_{12}+bE_{13}+cE_{23}\in R$和$Y=uE_{12}+vE_{13}\in N$，矩阵单位乘法$E_{ij}E_{k\ell}=\delta_{jk}E_{i\ell}$给出
\[
XY=0,\qquad YX=ucE_{13}\in N.
\]
加法与取负的封闭性由线性表达式得到，所以$N\lhd R$。同时$N^2=0$，因此$I$对$N$的左右乘法均吸收为零，又是加法子群，故$I\lhd N$。

但是$E_{12}\in I$、$E_{23}\in R$，而
\[
E_{12}E_{23}=E_{13}\notin I.
\]
所以$I$不是$R$的理想。此处$R$没有幺元，符合一般环不预设幺元的约定。

\end{proof}

\paragraph{复习提示}

验证$I\lhd N$只检查$N$中元素，验证$I\lhd R$则必须检查整个$R$中的元素。



\subsection{习题20：非单位属于某个极大理想}

\label{ex:3:20}

\paragraph{题目}

设$R$为交换环、$a$不可逆，证明有极大理想包含$a$。

\paragraph{考点与思路}

从真主理想$(a)$出发，以佐恩引理选极大元，并检查它在全部真理想中也极大。

\paragraph{所用结论与原文依据}

\noteref{c3prime}；\noteref{c3principal}；\noteref{c3maxexists}。原解位置：\path{chapters/Part 2/chapter3-exercises.tex}，第213行起。本题在现有习题解答的思路内补齐推导。

% 原题和原解来源同上；新增中间推导为本册补写。

\begin{proof}

这里$R$为交换幺环，$a$为非单位。若$(a)=R$，则$1=ra$对某个$r\in R$成立；交换性使$r$为$a$的双侧逆，矛盾。因此$(a)\ne R$。

令$\mathcal S=\{I\lhd R:(a)\subseteq I\ne R\}$，以包含关系赋序。它含$(a)$。设$\mathcal C\subseteq\mathcal S$为非空链，令$U=\bigcup_{I\in\mathcal C}I$。任取$x,y\in U$，可在链中找一个同时含二者的理想，故$x-y\in U$；对$r\in R$，$rx\in U$。所以$U$是包含$(a)$的理想。

若$U=R$，则$1\in U$，从而链中某个$I$含$1$，使$I=R$，矛盾。因此$U\in\mathcal S$，是该链的上界。空链可取$(a)$为上界。由佐恩引理，$\mathcal S$有极大元$M$。

若$M\subsetneq J\subsetneq R$且$J$为理想，则$J$也含$(a)$，故$J\in\mathcal S$，违背$M$的极大性。因此$M$为$R$的极大理想，并含$a$。

\end{proof}

\paragraph{复习提示}

链的并仍是真理想，关键在于一旦并含$1$，某个链成员就已含$1$。



\subsection{习题21：极大理想之交的单位刻画}

\label{ex:3:21}

\paragraph{题目}

设$J(R)$为交换环$R$的全部极大理想之交。证明$a\in J(R)$当且仅当$1-ax$对所有$x\in R$均可逆。

\paragraph{考点与思路}

按原解分别把非单位送入极大理想，以及在商域中取$a$的逆。

\paragraph{所用结论与原文依据}

\noteref{c3prime}；\exref{3}{20}。原解位置：\path{chapters/Part 2/chapter3-exercises.tex}，第222行起。本题在现有习题解答的思路内补齐推导。

% 原题和原解来源同上；新增中间推导为本册补写。

\begin{proof}

设$R$为交换幺环，记$J(R)=\bigcap_{M\text{ 极大}}M$。若$R$为零环，按空交为$R$的约定，$J(R)=R$；唯一元素也是单位，故等价性成立。以下设$R\ne0$。

先设$a\in J(R)$。任取$x\in R$。若$1-ax$不是单位，由\exref{3}{20}，它属于某个极大理想$M$。又$a\in J(R)\subseteq M$，理想的吸收性说明$ax\in M$。于是
\[
1=(1-ax)+ax\in M,
\]
与$M$为真理想矛盾。因此$1-ax$为单位。

反过来，设每个$1-ax$均可逆。如果$a\notin J(R)$，按交集的定义存在极大理想$M$使$a\notin M$。商环$R/M$是域，非零元素$a+M$有逆。取该逆的一个代表元$x\in R$，则
\[
(a+M)(x+M)=1+M,\qquad 1-ax\in M.
\]
但真理想不含单位：若单位$u\in M$，则$u^{-1}u=1\in M$。因此$1-ax$不可逆，与假设矛盾。故$a\in J(R)$。

\end{proof}

\paragraph{复习提示}

对所有$x$的量词不能删去；反向证明要在商域中取逆并提升代表元。



\subsection{习题22：强拟正则元全体与雅各布森根}

\label{ex:3:22}

\paragraph{题目}

若$a$满足习题21的条件，称其为强拟正则元。证明全部强拟正则元构成$J(R)$。

\paragraph{考点与思路}

把强拟正则的定义逐元素代入上一题的等价性，再说明所得集合为理想。

\paragraph{所用结论与原文依据}

\noteref{ideal}；\exref{3}{21}。原解位置：\path{chapters/Part 2/chapter3-exercises.tex}，第231行起。本题在现有习题解答的思路内补齐推导。

% 原题和原解来源同上；新增中间推导为本册补写。

\begin{proof}

在交换幺环$R$中，称$a$为强拟正则元，是指对每个$x\in R$，$1-ax$都是单位。记全部这类元素的集合为$Q$。由\exref{3}{21}，对任意$a\in R$都有
\[
a\in Q\ \Longleftrightarrow\ (\forall x\in R)\ 1-ax\in R^\times
\ \Longleftrightarrow\ a\in J(R).
\]
因此$Q=J(R)$。$J(R)$作为极大理想的交是理想：若$a,b$属于交，则$a-b$属于每个极大理想；对$r\in R$，$ra$也属于每个极大理想。所以强拟正则元全体也构成理想。零环按上一题的空交约定处理，结论相同。

\end{proof}

\paragraph{复习提示}

此题中“强拟正则”采用对每个乘子$x$均成立的定义。



\subsection{习题23：有限交换环的素理想必极大}

\label{ex:3:23}

\paragraph{题目}

证明有限交换环的每个素理想都是极大理想。

\paragraph{考点与思路}

把素理想的商环化为有限整环，再用有限整环为域。

\paragraph{所用结论与原文依据}

\noteref{c3prime}；\exref{3}{18}。原解位置：\path{chapters/Part 2/chapter3-exercises.tex}，第242行起。本题在现有习题解答的思路内补齐推导。

% 原题和原解来源同上；新增中间推导为本册补写。

\begin{proof}

设$P$为有限交换幺环$R$的素理想。素理想为真理想，故$R/P\ne0$；素性给出
\[
(a+P)(b+P)=0+P\ \Longrightarrow\ ab\in P
\ \Longrightarrow\ a\in P\text{ 或 }b\in P.
\]
所以$R/P$无零因子，是整环。它由有限集合$R$取商得到，因而仍有限。由\exref{3}{18}，$R/P$是域。根据极大理想与商域的等价判据，$P$极大。

\end{proof}

\paragraph{复习提示}

先确认商环非零，才能使用有限整环为域。



\subsection{习题24：整数子群定理与循环群分类的等价性}

\label{ex:3:24}

\paragraph{题目}

证明第二章第2节的定理4和定理5等价。

\paragraph{考点与思路}

按原解使用从整数群到循环群的指数同态，再用无限循环群返回整数子群。

\paragraph{所用结论与原文依据}

\noteref{cyclic}；\noteref{firstiso}；\noteref{correspondence}。原解位置：\path{chapters/Part 2/chapter3-exercises.tex}，第249行起。本题在现有习题解答的思路内补齐推导。

% 原题和原解来源同上；新增中间推导为本册补写。

\begin{proof}

本题所说的旧编号“第二章第2节定理4、5”，在现有源码中分别是整数子群命题（\path{chapters/Part 2/chapter2.tex}第210行起）及紧随的循环群分类定理（第228行起）。这里按内容比较，不把旧编号当成当前统一编号。

设已知整数子群定理：$\Z$的每个加法子群唯一形如$m\Z$，$m\ge0$。取$G=\langle g\rangle$，定义满同态
\[
\varphi:\Z\longrightarrow G,\qquad k\longmapsto g^k.
\]
其核为$n\Z$。若$n=0$，$\varphi$单射，故$G\cong\Z$。若$n>0$，同态基本定理给出$G\cong\Z/n\Z$，故$|G|=n$。

对任意子群$H\le G$，完全反像$\varphi^{-1}(H)$是$\Z$的子群，故为$m\Z$。满射性给出
\[
H=\varphi(m\Z)=\langle g^m\rangle.
\]
无限情形中，这已说明所有子群都为循环群。有限情形中，$n\Z=\ker\varphi\subseteq m\Z$，所以$m\mid n$且$m>0$。元素$e,g^m,\ldots,g^{(n/m-1)m}$恰为不同的全部元素：两项相等当且仅当$n\mid m(i-j)$，也就是$n/m\mid i-j$。因此$|H|=n/m$。给定$s\mid n$，只能取$m=n/s$，得到唯一的$s$阶子群$\langle g^{n/s}\rangle$。这证明循环群分类。

反过来，假设循环群分类成立。$(\Z,+)=\langle1\rangle$是无限循环群，每个子群都是循环群，故为$\langle m\rangle=m\Z$。取$m\ge0$；零子群只对应$m=0$，非零子群中$m$为最小正元素，故参数唯一。又
\[
n\Z\supseteq m\Z\ \Longleftrightarrow\ m\in n\Z
\ \Longleftrightarrow\ n\mid m.
\]
所以整数子群定理也得到证明。

\end{proof}

\paragraph{复习提示}

同态定理与子群对应是两种表述之间的共同前置工具；旧编号须按具体内容定位。



\subsection{习题25：单位根群的满而非单自同态}

\label{ex:3:25}

\paragraph{题目}

设$G$为复数域中全部单位根构成的乘法群。给出一个满而非单的自同态。

\paragraph{考点与思路}

保留平方映射，取平方根证明满射，用$1$和$-1$证明不单。

\paragraph{所用结论与原文依据}

\noteref{hom}。原解位置：\path{chapters/Part 2/chapter3-exercises.tex}，第260行起。本题在现有习题解答的思路内补齐推导。

% 原题和原解来源同上；新增中间推导为本册补写。

\begin{proof}

记$G=\{z\in\C:z^m=1\text{ 对某个 }m\ge1\}$，定义$\varphi:G\to G$为$\varphi(z)=z^2$。
若$z^m=1$，则$(z^2)^m=1$，所以$\varphi$确实映入$G$。因复数乘法交换，
\[
\varphi(zw)=(zw)^2=z^2w^2=\varphi(z)\varphi(w),
\]
它是自同态。

任取$w\in G$，取$m\ge1$使$w^m=1$。复数中存在$z$使$z^2=w$，于是$z^{2m}=w^m=1$，说明$z\in G$，且$\varphi(z)=w$。因此$\varphi$满射。
但$1\ne-1$均在$G$中，且$\varphi(1)=\varphi(-1)=1$，所以$\varphi$不单射。

\end{proof}

\paragraph{复习提示}

取出的平方根还须属于单位根群，而不只是属于复数域。



\subsection{习题26：若尔当同态的幂与极化恒等式}

\label{ex:3:26}

\paragraph{题目}

设$R,R'$为幺环，$\eta:R\to R'$为加法映射，$\eta(1)=1'$，且$\eta(aba)=\eta(a)\eta(b)\eta(a)$。说明其若尔当同态含义，证明$\eta(a^k)=\eta(a)^k$（$k\ge0$），并证明$\eta(ab+ba)$和$\eta(abc+cba)$的极化恒等式。

\paragraph{考点与思路}

按原解先代入$1$得到平方公式，再按奇偶归纳幂公式，最后展开两个和式。

\paragraph{所用结论与原文依据}

\noteref{ringlaws}；\noteref{unit}。原解位置：\path{chapters/Part 2/chapter3-exercises.tex}，第267行起。本题在现有习题解答的思路内补齐推导。

% 原题和原解来源同上；新增中间推导为本册补写。

\begin{proof}

为使后续题目的含义明确，本组26--29题采用保幺加法映射且满足$\eta(aba)=\eta(a)\eta(b)\eta(a)$的若尔当同态约定。它必保平方；这里只保留平方的较弱定义在特征$2$时不能直接替代三重乘积条件。

记$A=\eta(a)$。取$b=1$，由保幺性得
\[
\eta(a^2)=\eta(a)\eta(1)\eta(a)=A^2.
\]
又$\eta(a^0)=\eta(1)=1'$、$\eta(a)=A$。若$\eta(a^m)=A^m$，则用三重乘积公式得到
\[
\eta(a^{m+2})=\eta(aa^ma)=A\eta(a^m)A=A^{m+2}.
\]
分别以$m=0,1$为起点归纳，故对每个整数$k\ge0$，$\eta(a^k)=\eta(a)^k$。

由$(a+b)^2=a^2+ab+ba+b^2$，加法性及保平方性给出
\begin{align*}
\eta(a^2)+\eta(ab+ba)+\eta(b^2)
&=\eta((a+b)^2)\\
&=(\eta(a)+\eta(b))^2.
\end{align*}
在加法群中消去两个平方项，得到
\[
\eta(ab+ba)=\eta(a)\eta(b)+\eta(b)\eta(a).
\]
最后展开$(a+c)b(a+c)=aba+abc+cba+cbc$，得
\begin{align*}
\eta(aba)+\eta(abc+cba)+\eta(cbc)
&=(\eta(a)+\eta(c))\eta(b)(\eta(a)+\eta(c)).
\end{align*}
消去$\eta(aba)$及$\eta(cbc)$，得到
\[
\eta(abc+cba)=\eta(a)\eta(b)\eta(c)+\eta(c)\eta(b)\eta(a).
\]
两次展开只用加法消去，没有除以$2$，因而在特征$2$时也成立。

\end{proof}

\paragraph{复习提示}

三重乘积极化式来自三重乘积假设；不能只凭平方公式而在特征$2$中除以$2$。



\subsection{习题27：若尔当同态对一对元素的二择一性质}

\label{ex:3:27}

\paragraph{题目}

设$\eta:R\to R'$是到无零因子环的若尔当同态。证明对任意$a,b$，必有$\eta(ab)=\eta(a)\eta(b)$或$\eta(ab)=\eta(b)\eta(a)$。

\paragraph{考点与思路}

沿原解对三重乘积极化式代入$ab,b,a$，再按固定次序因式分解。

\paragraph{所用结论与原文依据}

\noteref{zerodiv}；\exref{3}{26}。原解位置：\path{chapters/Part 2/chapter3-exercises.tex}，第285行起。本题在现有习题解答的思路内补齐推导。

% 原题和原解来源同上；新增中间推导为本册补写。

\begin{proof}

沿用\exref{3}{26}的若尔当同态约定，记
\[
A=\eta(a),\quad B=\eta(b),\quad C=\eta(ab).
\]
在该题的极化式中令外侧变量为$ab,a$，中间变量为$b$，得到
\[
\eta(ab^2a+abab)=CBA+ABC.
\]
由三重乘积公式及平方公式，
\[
\eta(ab^2a)=A\eta(b^2)A=AB^2A,\qquad
\eta(abab)=\eta((ab)^2)=C^2.
\]
加法性给出$C^2+AB^2A=CBA+ABC$。保持因子顺序展开，有
\begin{align*}
(C-AB)(C-BA)
&=C^2-CBA-ABC+ABBA\\
&=C^2-CBA-ABC+AB^2A=0.
\end{align*}
目标环无零因子，故$C-AB=0$或$C-BA=0$。即对这对$a,b$，
\[
\eta(ab)=\eta(a)\eta(b)\quad\text{或}\quad\eta(ab)=\eta(b)\eta(a).
\]

\end{proof}

\paragraph{复习提示}

本题的二择一暂时依赖于$a,b$；还未证明所有元素对都选择同一乘法方向。



\subsection{习题28：逐对二择一推出全局同态或反同态}

\label{ex:3:28}

\paragraph{题目}

设$\eta:R\to R'$是除环之间的加法映射，$\eta(1)=1'$，且对每对$a,b$都有$\eta(ab)=\eta(a)\eta(b)$或$\eta(b)\eta(a)$。证明$\eta$为同态或反同态。

\paragraph{考点与思路}

沿原解对第一个变量、第二个变量分别使用一次“两子群之并”论证。

\paragraph{所用结论与原文依据}

\noteref{subgroup}；\noteref{c3anti}。原解位置：\path{chapters/Part 2/chapter3-exercises.tex}，第304行起。本题在现有习题解答的思路内补齐推导。

% 原题和原解来源同上；新增中间推导为本册补写。

\begin{proof}

先证所用群论事实：若群$G=H\cup K$且$H,K$均为子群，则$H=G$或$K=G$。否则二者都真，且互不包含，可取$h\in H\setminus K$、$k\in K\setminus H$。若$hk\in H$则$k=h^{-1}(hk)\in H$；若$hk\in K$则$h=(hk)k^{-1}\in K$，两种都矛盾。因此$hk$不属于并集，违背$G=H\cup K$。

固定$b\in R$，定义
\begin{align*}
H_b&=\{a:\eta(ab)=\eta(a)\eta(b)\},\\
K_b&=\{a:\eta(ab)=\eta(b)\eta(a)\}.
\end{align*}
二者都含$0$。若$a,c\in H_b$，加法性与分配律给出
\[
\eta((a-c)b)=\eta(ab)-\eta(cb)
=(\eta(a)-\eta(c))\eta(b)=\eta(a-c)\eta(b),
\]
所以$a-c\in H_b$。若$a,c\in K_b$，同样计算为
\[
\eta((a-c)b)=\eta(b)(\eta(a)-\eta(c))=\eta(b)\eta(a-c),
\]
故$K_b$也为加法子群。题设说明$R=H_b\cup K_b$，因而每个$b$都满足$H_b=R$或$K_b=R$。

令$B_+=\{b:H_b=R\}$、$B_-=\{b:K_b=R\}$。二者含$0$。若$b,c\in B_+$，则对所有$a$，
\[
\eta(a(b-c))=\eta(ab)-\eta(ac)=\eta(a)(\eta(b)-\eta(c))=\eta(a)\eta(b-c),
\]
所以$b-c\in B_+$。若$b,c\in B_-$，则
\[
\eta(a(b-c))=\eta(b)\eta(a)-\eta(c)\eta(a)=\eta(b-c)\eta(a),
\]
故$b-c\in B_-$。于是$B_+,B_-$都是加法子群，且$R=B_+\cup B_-$。再次使用所证事实，$B_+=R$或$B_-=R$。

第一种情形给出对所有$a,b$都有$\eta(ab)=\eta(a)\eta(b)$；第二种给出对所有$a,b$都有$\eta(ab)=\eta(b)\eta(a)$。结合既有加法性，$\eta$分别为同态或反同态。该证明实际不需要源、目标为除环。

\end{proof}

\paragraph{复习提示}

使用两次子群并，才能消除方向分别依赖于两个变量的可能。



\subsection{习题29：除环若尔当同态的乘法方向}

\label{ex:3:29}

\paragraph{题目}

证明除环到无零因子环的若尔当同态必为同态或反同态。

\paragraph{考点与思路}

将上一题的逐对结论送入双重子群论证，并核对目标环无需为除环。

\paragraph{所用结论与原文依据}

\exref{3}{26}；\exref{3}{27}；\exref{3}{28}。原解位置：\path{chapters/Part 2/chapter3-exercises.tex}，第318行起。本题在现有习题解答的思路内补齐推导。

% 原题和原解来源同上；新增中间推导为本册补写。

\begin{proof}

按\exref{3}{26}的约定，$\eta$保加法、幺元与三重乘积；目标环$R'$无零因子。因此\exref{3}{27}适用，对每对$a,b$，有
\[
\eta(ab)=\eta(a)\eta(b)\quad\text{或}\quad
\eta(ab)=\eta(b)\eta(a).
\]
\exref{3}{28}的证明只用了加法性、分配律和上述逐对二择一性质，并未对目标元素取逆。故虽本题只假设$R'$无零因子，仍可应用该证明，得到对所有$a,b$统一成立的同一方向。于是$\eta$为同态或反同态。

这里保留第26题的三重乘积假设是必要的定义区分。若把“若尔当同态”仅定义为保幺、加法且保平方，则在特征$2$时，本题结论一般不成立。取八元域$K$（可按\exref{3}{37}构造），定义
\[
T:K\longrightarrow\F_2,\qquad T(z)=z+z^2+z^4.
\]
因为$z^8=z$，有$T(z)^2=T(z)$，所以$T(z)$确在$\F_2$中。特征$2$的幂公式说明$T$保加法，且$T(1)=1$、$T(z^2)=T(z)^2$。它从八个元素映到两个元素，故不单射；而域之间的保幺同态必单射（非零核元素取逆会使$1$落入核），交换域的反同态也就是同态。因此$T$不是同态或反同态。

它也不满足本组保留的三重乘积条件：取$0\ne a\in\ker T$，令$b=a^{-2}$，则$aba=1$，却$T(a)T(b)T(a)=0\ne T(1)$。这说明不能把本组的三重乘积定义无条件换成仅保平方的定义。

\end{proof}

\paragraph{复习提示}

不得只从每对元素各有一种方向就直接断言全局方向；统一方向的理由在上一题的两个子群论证中。



\subsection{习题30：模素数幂环的零因子与素理想}

\label{ex:3:30}

\paragraph{题目}

设$R=\mathbb Z/(p^n)$，其中$p$为素数、$n>1$。确定零因子、素理想及相应商域。

\paragraph{考点与思路}

先用互素判别单位，再从商环理想对应列出全部理想，逐一检验素性。

\paragraph{所用结论与原文依据}

\noteref{c3correspondence}；\noteref{c3prime}；\noteref{zerodiv}。原解位置：\path{chapters/Part 2/chapter3-exercises.tex}，第325行起。本题在现有习题解答的思路内补齐推导。

% 原题和原解来源同上；新增中间推导为本册补写。

\begin{proof}

记$\bar a=a+(p^n)$。若$p\nmid a$，则$\gcd(a,p^n)=1$，存在整数$s,t$使$sa+tp^n=1$，故$\bar a$为单位。若$p\mid a$，则其模$p$像为零，不可能可逆；否则模$p$也会有逆。

零因子要求非零。若$\bar a\ne0$且$p\mid a$，可写$a=p^su$，其中$1\le s<n$且$p\nmid u$，于是
\[
\bar a\,\overline{p^{n-s}}=0,\qquad \overline{p^{n-s}}\ne0.
\]
故非零的$p$倍数均为零因子。其余非零元素是单位，不可能消去非零元素。因此零因子集恰为
\[
(p)/(p^n)\setminus\{0+(p^n)\}.
\]
注意非单位集还含$0$，不能不加区别地与零因子集等同。

由$\Z\to R$的理想对应，$R$的理想来自包含$(p^n)$的整数理想$(d)$，即$d\mid p^n$。所以全部理想为$(p^k)/(p^n)$，$0\le k\le n$。$k=0$时为整个环，不素；若$2\le k\le n$，则$\bar p\,\overline{p^{k-1}}\in(p^k)/(p^n)$，但两个因子都不在其中，所以不素。

模$p$的满同态核为$P=(p)/(p^n)$，故
\[
R/P\cong\Z/(p)=\F_p.
\]
这是域，因而$P$是唯一素理想，也是唯一极大理想；相应商域为$\F_p$。

\end{proof}

\paragraph{复习提示}

$0$是非单位却按本笔记约定不是零因子；列出素理想时须排除整个环。



\subsection{习题31：模素数幂系数多项式的单位与零因子}

\label{ex:3:31}

\paragraph{题目}

设$R=\mathbb Z/(p^n)$，确定$R[x]$的单位和零因子。

\paragraph{考点与思路}

保留模$p$与系数幂零的判别；为避免跨题缺口，在题内证明单位的必要条件，并展开原解的最低赋值系数论证。

\paragraph{所用结论与原文依据}

\noteref{c3polyquot}；\noteref{c3degree}；\exref{3}{30}。原解位置：\path{chapters/Part 2/chapter3-exercises.tex}，第336行起。本题在现有习题解答的思路内补齐推导。

% 原题和原解来源同上；新增中间推导为本册补写。

\begin{proof}

设$n\ge1$、$f=\sum_{i=0}^m a_ix^i\in(\Z/(p^n))[x]$。先求单位。若$fg=1$，比较常数项知$a_0$为单位。逐系数模$p$约化后，在整环$\F_p[x]$中有$\bar f\bar g=1$。次数可加，所以$\bar f$为非零常数；因此$p\nmid a_0$，而所有$i>0$的$a_i$均被$p$整除。

反过来，若满足这些系数条件，$a_0$可逆，令$h=a_0^{-1}(f-a_0)$。$h$的每个系数被$p$整除，所以$h^n=0$。于是
\[
f^{-1}=a_0^{-1}\sum_{j=0}^{n-1}(-h)^j,
\]
左右相乘均由有限等比和得到$1$。故
\[
f\text{ 可逆}\ \Longleftrightarrow\ p\nmid a_0\text{ 且 }p\mid a_i\ (i>0).
\]

再求非零零因子。若$f\ne0$且全部系数均被$p$整除，则常数$\overline{p^{n-1}}\ne0$且$\overline{p^{n-1}}f=0$，故$f$为零因子（$n=1$时不存在这样的非零$f$）。

反过来，设$f$有不被$p$整除的系数，取最小的这类下标$j$。对剩余类$r$，令$v(r)$为最大的$s\in\{0,\ldots,n\}$，使$r$被$p^s$整除，其中$v(0)=n$。任取非零$g=\sum b_kx^k$，令$s=\min_kv(b_k)<n$，再取达到该最小值的最小下标$k$。

在$fg$的$x^{j+k}$项中，$a_jb_k$的赋值恰为$s$。其余乘积$a_ib_\ell$满足$i+\ell=j+k$。若$i<j$，则$v(a_i)\ge1$而$v(b_\ell)\ge s$；若$i>j$，则$\ell<k$，由$k$的选择有$v(b_\ell)\ge s+1$。故所有其他项都被$p^{s+1}$整除，$a_jb_k$却不被$p^{s+1}$整除。总和不能为零模$p^n$，因$s<n$。所以$fg\ne0$。

因此非零零因子恰为全部系数被$p$整除的非零多项式。

\end{proof}

\paragraph{复习提示}

以最低次单位系数配合最小赋值的最低次系数，才能保证选中的乘积不会被其他项抵消。



\subsection{习题32：直和环的理想分解及整数商环的根与单位群}

\label{ex:3:32}

\paragraph{题目}

设$R=R_1\oplus\cdots\oplus R_r$，$N\lhd R$。证明$N=\bigoplus_i(N\cap R_i)$，并求$\mathbb Z/(n)$的最大幂零理想和单位群。

\paragraph{考点与思路}

用各分量的单位元投影理想，再按中国剩余分解逐个研究模素数幂因子。

\paragraph{所用结论与原文依据}

\noteref{c3direct}；\noteref{c3crt}；\noteref{c3nil}；\exref{3}{30}。原解位置：\path{chapters/Part 2/chapter3-exercises.tex}，第345行起。本题在现有习题解答的思路内补齐推导。

% 原题和原解来源同上；新增中间推导为本册补写。

\begin{proof}

\begin{enumerate}
\item 设每个$R_i$为幺环，并将它等同于$R=\bigoplus\limits_{i=1}^rR_i$中仅第$i$分量非零的理想。令$e_i$为该分量的单位元，其他分量为零。对$x=(x_1,\ldots,x_r)\in N$，理想吸收性给出
\[
e_ix=(0,\ldots,x_i,\ldots,0)\in N\cap R_i,\qquad
x=\sum_{i=1}^r e_ix.
\]
所以$N\subseteq\sum_i(N\cap R_i)$；反向包含来自每个交均含于$N$。不同坐标的表示唯一，因此
\[
N=\bigoplus\limits_{i=1}^r(N\cap R_i).
\]
若不要求各分量有幺元，取两个零乘法环$\Z$的直和，斜子群$\{(m,m):m\in\Z\}$就是反例。
\item 设$n\ge2$，写$n=\prod_i p_i^{a_i}$，令$d=\prod_i p_i$。中国剩余定理给出$\Z/(n)\cong\bigoplus_i\Z/(p_i^{a_i})$。在第$i$个因子中，$p_i$的倍数取$a_i$次幂即为零；若不被$p_i$整除，则是单位，非零环中的单位不可能幂零。因此幂零元恰为每个分量中被$p_i$整除的元。整数$a$的剩余类满足此条件当且仅当$d\mid a$。

故幂零元理想为$N=(d)/(n)$。它本身是幂零理想：取$A=\max_i a_i$，因$n\mid d^A$，有$N^A=0$。任何幂零理想的每个元素都幂零，故包含于$N$。因此$N$确为最大的幂零理想。
\item 分量乘积为$1$等价于每个分量乘积为$1$，故
\[
(\Z/(n))^\times\cong\prod_i(\Z/(p_i^{a_i}))^\times.
\]
在模$p^a$环中，不是单位的元素恰为$p$的倍数，共$p^{a-1}$个，所以单位有$p^a-p^{a-1}$个。因此
\[
|(\Z/(n))^\times|=\prod_i(p_i^{a_i}-p_i^{a_i-1})
=n\prod_{p\mid n}\left(1-\frac1p\right).
\]
以上直积分解给出单位群；并未无依据地把每个因子再断言为循环群。$n=1$时零环的最大幂零理想为零理想，单位群为单元素群。
\end{enumerate}

\end{proof}

\paragraph{复习提示}

一般幂零元理想未必幂零；本题额外证明$N^A=0$，才得到最大幂零理想。



\subsection{习题33：三次多项式在六元环中的六个根}

\label{ex:3:33}

\paragraph{题目}

证明$x^3-x$在$\mathbb Z/(6)$中有六个根。

\paragraph{考点与思路}

沿用连续三个整数的整除性，并说明取代表元不影响代入值。

\paragraph{所用结论与原文依据}

\noteref{c3roots}；\noteref{c3poly}。原解位置：\path{chapters/Part 2/chapter3-exercises.tex}，第364行起。本题在现有习题解答的思路内补齐推导。

% 原题和原解来源同上；新增中间推导为本册补写。

\begin{proof}

对任意整数$a$，
\[
a^3-a=(a-1)a(a+1).
\]
三个连续整数中至少一个为偶数，且恰有一个被$3$整除，所以$2\mid a^3-a$且$3\mid a^3-a$。因$2,3$互素，$6\mid a^3-a$。

因此对每个$\bar a\in\Z/(6)$，都有$\bar a^3-\bar a=0$；模$6$的代入只依赖于剩余类，因为自然同态保加法和乘法。于是全部六个元素$\bar0,\bar1,\ldots,\bar5$都是根，且环中没有其他元素，恰有六个根。

\end{proof}

\paragraph{复习提示}

次数限制根数的结论要求系数环为整环；$\Z/(6)$有零因子。



\subsection{习题34：四元数中的二次方程有无穷多个根}

\label{ex:3:34}

\paragraph{题目}

证明$x^2+1$在四元数体$H$中有无穷多个根。

\paragraph{考点与思路}

按原习题解展开纯虚四元数的平方，给出明确的无限族来证明互异性。

\paragraph{所用结论与原文依据}

\noteref{quaternion}；\noteref{c3qroots}。原解位置：\path{chapters/Part 2/chapter3-exercises.tex}，第371行起。本题在现有习题解答的思路内补齐推导。

% 原题和原解来源同上；新增中间推导为本册补写。

\begin{proof}

取纯虚四元数$q=bI+cJ+dK$。由乘法表，$IJ+JI=IK+KI=JK+KJ=0$，所以
\begin{align*}
q^2&=b^2I^2+c^2J^2+d^2K^2
+bc(IJ+JI)+bd(IK+KI)+cd(JK+KJ)\\
&=-(b^2+c^2+d^2).
\end{align*}
因此每个满足$b^2+c^2+d^2=1$的三元组都给出$x^2+1$的根。

例如对每个$t\in[0,1]$，令
\[
q_t=\frac{1-t^2}{1+t^2}I-\frac{2t}{1+t^2}K.
\]
其系数平方和为$((1-t^2)^2+4t^2)/(1+t^2)^2=1$，故$q_t^2=-1$。若$0\le s<t\le1$，则
\[
\frac{1-s^2}{1+s^2}-\frac{1-t^2}{1+t^2}
=\frac{2(t^2-s^2)}{(1+s^2)(1+t^2)}>0,
\]
所以$q_s\ne q_t$。因此有无穷多个根。这个族也等于$(1+tJ)I(1+tJ)^{-1}$，与原笔记的共轭构造一致。

\end{proof}

\paragraph{复习提示}

即使除环没有零因子，乘法不交换时也不能直接套用整环上的根数定理。



\subsection{习题35：两个平方根之和的极小多项式}

\label{ex:3:35}

\paragraph{题目}

设$u=\sqrt2+\sqrt3$，求其在$\mathbb Q$上的极小多项式。

\paragraph{考点与思路}

保留原解先求四次关系、再恢复两个平方根的路线；展开次数恰为四的理由。

\paragraph{所用结论与原文依据}

\noteref{c3minimal}；\noteref{c3division}。原解位置：\path{chapters/Part 2/chapter3-exercises.tex}，第382行起。本题在现有习题解答的思路内补齐推导。

% 原题和原解来源同上；新增中间推导为本册补写。

\begin{proof}

令$u=\sqrt2+\sqrt3$。有
\[
u^2=5+2\sqrt6,\qquad (u^2-5)^2=24,
\]
故$u^4-10u^2+1=0$。候选多项式为$f=x^4-10x^2+1$。

先验证$\sqrt3\notin\Q(\sqrt2)$。若$\sqrt3=r+s\sqrt2$，$r,s\in\Q$，平方并利用$\sqrt2$无理，得
\[
3=r^2+2s^2,
\qquad 2rs=0.
\]
若$s=0$，则$r^2=3$；若$r=0$，则$s^2=3/2$。二者均不可能为有理数的平方，因为有理数平方的任一素因子指数均为偶数，而$3$在这两个数中的指数均为奇数。于是$1,\sqrt2,\sqrt3,\sqrt6$在$\Q$上线性无关：若
\[
(a+b\sqrt2)+(c+d\sqrt2)\sqrt3=0,
\]
且$c+d\sqrt2\ne0$，可在域$\Q(\sqrt2)$中相除得到$\sqrt3\in\Q(\sqrt2)$，矛盾；故$c=d=0$，继而$a=b=0$。

又$(\sqrt3+\sqrt2)(\sqrt3-\sqrt2)=1$，所以$u^{-1}=\sqrt3-\sqrt2$，从而
\[
\sqrt3=\frac{u+u^{-1}}2,\qquad
\sqrt2=\frac{u-u^{-1}}2.
\]
由$f(u)=0$还可得$u^{-1}=10u-u^3\in\Q[u]$，故$\Q[u]$含上述四个线性无关元素。另一方面，$u^4=10u^2-1$使每个$u$的高次幂都化为$1,u,u^2,u^3$的有理线性组合。

若存在次数$d<4$的非零多项式$g$使$g(u)=0$，用其非零首项系数相除后，可以将$u^d$及全部更高幂降为$1,u,\ldots,u^{d-1}$的线性组合。这会使$\Q[u]$的维数至多$d<4$，与含四个线性无关元素矛盾。因此$f$是次数最小的首一关系，即极小多项式。

\end{proof}

\paragraph{复习提示}

找到一个四次关系只证明极小次数不超过四，还必须排除低次关系。



\subsection{习题36：三次代数关系下的乘法和求逆}

\label{ex:3:36}

\paragraph{题目}

设$u$属于包含$\Q$的域，且$u^3-u+1=0$。将$(5u^2+3u-1)(2u^2-2u+6)$及$(3u^2-u+2)^{-1}$写成次数不超过$2$的多项式。补明原解中除以$7$所用的特征$0$条件。

\paragraph{考点与思路}

按原解将$u^3,u^4$化为低次幂，并逐项展开验证逆元。

\paragraph{所用结论与原文依据}

\noteref{c3poly}；\noteref{c3minimal}。原解位置：\path{chapters/Part 2/chapter3-exercises.tex}，第394行起。本题在现有习题解答的思路内补齐推导。

% 原题和原解来源同上；新增中间推导为本册补写。

\begin{proof}

这里在特征为$0$的包含$\Q$的域中计算；这使$7$可逆。由$u^3-u+1=0$，得$u^3=u-1$、$u^4=u^2-u$。

首先展开乘积：
\begin{align*}
(5u^2+3u-1)(2u^2-2u+6)
&=10u^4-4u^3+22u^2+20u-6\\
&=10(u^2-u)-4(u-1)+22u^2+20u-6\\
&=32u^2+6u-2.
\end{align*}
对求逆，按原解取候选$(-2u^2-u+3)/7$。直接展开，
\begin{align*}
(3u^2-u+2)(-2u^2-u+3)
&=-6u^4-u^3+6u^2-5u+6\\
&=-6(u^2-u)-(u-1)+6u^2-5u+6\\
&=7.
\end{align*}
这些因子都是$u$的有理系数多项式，彼此交换，所以反向乘积也为$7$。除以$7$得到
\[
(3u^2-u+2)^{-1}=\frac{-2u^2-u+3}{7}.
\]

\end{proof}

\paragraph{复习提示}

求逆的结果必须实际相乘核对；原式若在任意特征下讨论，特征$7$时不能使用这一除法。



\subsection{习题37：八元域的商环构造}

\label{ex:3:37}

\paragraph{题目}

证明$\mathbb F_2[x]/(x^3+x^2+1)$是八元域。

\paragraph{考点与思路}

先在两个元素上检验三次式无根，再以带余除法数出商环的全部代表元。

\paragraph{所用结论与原文依据}

\noteref{c3irred}；\noteref{c3division}。原解位置：\path{chapters/Part 2/chapter3-exercises.tex}，第409行起。本题在现有习题解答的思路内补齐推导。

% 原题和原解来源同上；新增中间推导为本册补写。

\begin{proof}

令$h=x^3+x^2+1\in\F_2[x]$。有$h(0)=1$、$h(1)=1+1+1=1$，所以$h$在$\F_2$中无根。
若域上的三次多项式可约，则可分成两个正次数因子，次数和为$3$，必有一个因子的次数为$1$；一次因子的根也是原多项式的根。因此三次多项式无根即不可约，$h$不可约。

由不可约多项式的商环判据，$\F_2[x]/(h)$为域。用$h$作带余除法，每个剩余类有代表元$a+bx+cx^2$，$a,b,c\in\F_2$。若两个这样的代表元同余，其差是次数小于$3$的$h$的倍式，只能为零，所以表示唯一。三系数独立各取两值，因此恰有$2^3=8$个元素。

\end{proof}

\paragraph{复习提示}

无根等价于不可约仅适用于二次或三次式，不能推广到任意次数。



\subsection{习题38：五元域上的全部首一不可约二次式}

\label{ex:3:38}

\paragraph{题目}

列出$\mathbb F_5[x]$中全部首项系数为$1$的二次不可约多项式。

\paragraph{考点与思路}

保留判别式路线，以表格逐一列出$b$与两种非平方判别式对应的$c$。

\paragraph{所用结论与原文依据}

\noteref{c3irred}；\noteref{c3roots}。原解位置：\path{chapters/Part 2/chapter3-exercises.tex}，第416行起。本题在现有习题解答的思路内补齐推导。

% 原题和原解来源同上；新增中间推导为本册补写。

\begin{proof}

考虑$f=x^2+bx+c\in\F_5[x]$。域上的二次式可约当且仅当有根，而
\[
4f(x)=(2x+b)^2-(b^2-4c).
\]
因$2$可逆，映射$x\mapsto2x+b$为双射，所以$f$有根当且仅当$\Delta=b^2-4c$是平方（含$0$）。$\F_5$的平方为$0,1,4$，非平方恰为$2,3$。

模$5$下$4=-1$，故$\Delta=b^2+c$，即$c=\Delta-b^2$。计算如下：
\begin{center}\begin{tabular}{c|c|cc}
$b$&$b^2$&$\Delta=2$时的$c$&$\Delta=3$时的$c$\\\hline
0&0&2&3\\1&1&1&2\\2&4&3&4\\3&4&3&4\\4&1&1&2
\end{tabular}\end{center}
因此全部所求多项式为
\begin{align*}
&x^2+2,\quad x^2+3,\\
&x^2+x+1,\quad x^2+x+2,\\
&x^2+2x+3,\quad x^2+2x+4,\\
&x^2+3x+3,\quad x^2+3x+4,\\
&x^2+4x+1,\quad x^2+4x+2.
\end{align*}
表格遍历了全部五个$b$及全部非平方判别式，每对$b,c$唯一决定一个多项式，故无遗漏或重复。

\end{proof}

\paragraph{复习提示}

判别式为零也可约；必须取非平方$2,3$，而不是仅要求判别式非零。



\subsection{习题39：二十五元域的构造}

\label{ex:3:39}

\paragraph{题目}

构造一个有$25$个元素的域。

\paragraph{考点与思路}

直接使用上一题的不可约二次式，并给出商环中的运算和唯一表示。

\paragraph{所用结论与原文依据}

\noteref{c3irred}；\noteref{c3division}；\exref{3}{38}。原解位置：\path{chapters/Part 2/chapter3-exercises.tex}，第427行起。本题在现有习题解答的思路内补齐推导。

% 原题和原解来源同上；新增中间推导为本册补写。

\begin{proof}

由\exref{3}{38}，$x^2+2$在$\F_5[x]$中不可约。令
\[
F=\F_5[x]/(x^2+2),\qquad \alpha=x+(x^2+2).
\]
不可约判据保证$F$是域。带余除法说明每个元素唯一写为$a+b\alpha$，其中$a,b\in\F_5$，因此共有$5^2=25$个元素。

其运算由$\alpha^2=-2=3$确定：
\begin{align*}
(a+b\alpha)+(c+d\alpha)&=(a+c)+(b+d)\alpha,\\
(a+b\alpha)(c+d\alpha)&=(ac+3bd)+(ad+bc)\alpha.
\end{align*}
为核对非零元可逆，若$(a,b)\ne(0,0)$，则
\[
(a+b\alpha)(a-b\alpha)=a^2+2b^2\ne0.
\]
若$b\ne0$而右端为零，则$(a/b)^2=-2=3$，与$3$非平方矛盾；若$b=0$，则$a\ne0$也使右端非零。因此逆元为$(a-b\alpha)/(a^2+2b^2)$。

\end{proof}

\paragraph{复习提示}

不能用$\Z/(25)$代替该域；后者有非零零因子$\bar5$。



\subsection{习题40：整数多项式环中的一个极大理想}

\label{ex:3:40}

\paragraph{题目}

证明$(3,x^3+2x^2+2x-1)$是$\mathbb Z[x]$的极大理想。

\paragraph{考点与思路}

保留模$3$化为三次不可约式的路线；把复合同态的核完整写出。

\paragraph{所用结论与原文依据}

\noteref{c3irred}；\noteref{c3polyquot}；\noteref{c3prime}。原解位置：\path{chapters/Part 2/chapter3-exercises.tex}，第438行起。本题在现有习题解答的思路内补齐推导。

% 原题和原解来源同上；新增中间推导为本册补写。

\begin{proof}

记$f=x^3+2x^2+2x-1\in\Z[x]$、$h=x^3-x^2-x-1\in\F_3[x]$。系数模$3$约化把$f$送为$h$。
逐一代入$\F_3$中的元素，有
\[
h(0)=2,\qquad h(1)=1,\qquad h(2)=1.
\]
三次式无根则不可约，故$K=\F_3[x]/(h)$是域。

考虑满同态$\psi:\Z[x]\to K$，先约化系数再取模$(h)$。若$g\in(3,f)$，则$\psi(g)=0$。反过来，若$\psi(g)=0$，则$\bar g=h\bar q$对某个$\bar q\in\F_3[x]$成立。选一个整数系数提升$q$，有$\overline{g-fq}=0$，故$g-fq$的每个系数均被$3$整除，写成$g=fq+3r$。因此
\[
\ker\psi=(3,f).
\]
由同态基本定理，$\Z[x]/(3,f)\cong K$是域，所以$(3,f)$为极大理想。

\end{proof}

\paragraph{复习提示}

商域同构要说明核恰为$(3,f)$，不能只证明两个生成元映为零。



\subsection{习题41：模三约化诱导的商环同构}

\label{ex:3:41}

\paragraph{题目}

令$I=(3,x^3+2x^2+2x-1)$。证明$\mathbb Z[x]/I\cong\mathbb F_3[x]/(x^3-x^2-x-1)$，并推出$I$极大。

\paragraph{考点与思路}

沿原解系数约化的核与像，明确诱导映射并使用上一题的核计算。

\paragraph{所用结论与原文依据}

\noteref{c3polyquot}；\noteref{c3prime}；\exref{3}{40}。原解位置：\path{chapters/Part 2/chapter3-exercises.tex}，第449行起。本题在现有习题解答的思路内补齐推导。

% 原题和原解来源同上；新增中间推导为本册补写。

\begin{proof}

记$I=(3,f)$，其中$f=x^3+2x^2+2x-1$，并记$h=x^3-x^2-x-1\in\F_3[x]$。定义
\[
\Phi:\Z[x]/I\longrightarrow\F_3[x]/(h),\qquad
g+I\longmapsto\bar g+(h),
\]
其中$\bar g$表示逐系数模$3$约化。

由\exref{3}{40}的完整核计算，$\bar g\in(h)$当且仅当$g\in I$。因此若$g-g'\in I$，则$\bar g-\bar g'\in(h)$，$\Phi$良定义；若$\Phi(g+I)=0$，则$g\in I$，故$\Phi$单射。任一$\F_3$系数多项式都有整数系数提升，所以$\Phi$满射。系数约化与取商均保加法、乘法，故$\Phi$是环同构。

上一题已逐个检查$h$的三个可能根，证明它不可约。因此右侧为域，左侧也为域，故$I$极大。

\end{proof}

\paragraph{复习提示}

本题与上一题共享核计算和不可约性检验，但仍应明确写出所要求的同构映射。



\subsection{习题42：多项式商环无零因子的判据}

\label{ex:3:42}

\paragraph{题目}

设$F$为域且$\deg f>0$。证明$F[x]/(f)$无零因子当且仅当$f$不可约。

\paragraph{考点与思路}

保留不可约时商环为域、可约时产生两个非零零乘积的双向证明。

\paragraph{所用结论与原文依据}

\noteref{c3irred}；\noteref{c3degree}。原解位置：\path{chapters/Part 2/chapter3-exercises.tex}，第456行起。本题在现有习题解答的思路内补齐推导。

% 原题和原解来源同上；新增中间推导为本册补写。

\begin{proof}

设$\deg f=d>0$。若$f$不可约，则$(f)$为极大理想，故$F[x]/(f)$是域，特别地无零因子。此处$f$次数为正，保证$(f)\ne F[x]$，商环非零。

反过来，设$f$可约。可写$f=gh$，其中$\deg g,\deg h>0$。在域上非零乘积次数可加，所以
\[
0<\deg g<d,\qquad 0<\deg h<d.
\]
若$g\in(f)$且$g\ne0$，则$g=qf$且$q\ne0$，使$\deg g=\deg q+d\ge d$，矛盾。因此$g+(f)\ne0$；同理$h+(f)\ne0$。但
\[
(g+(f))(h+(f))=f+(f)=0+(f),
\]
商环有非零零因子。因此商环无零因子时$f$必须不可约。

\end{proof}

\paragraph{复习提示}

两个因子在母环非零不够，还必须证明它们取模后仍非零。



\subsection{习题43：幂零多项式的每个系数均幂零}

\label{ex:3:43}

\paragraph{题目}

设$R$为交换环。若$f\in R[x]$幂零，证明$f$的每个系数均幂零。

\paragraph{考点与思路}

按原解先取最高系数，再模去其生成的幂零理想，以次数归纳处理其余系数。

\paragraph{所用结论与原文依据}

\noteref{c3polyquot}；\noteref{c3nil}；\noteref{c3nilideal}。原解位置：\path{chapters/Part 2/chapter3-exercises.tex}，第463行起。本题在现有习题解答的思路内补齐推导。

% 原题和原解来源同上；新增中间推导为本册补写。

\begin{proof}

对$f$的次数归纳。若$f=0$，所有系数为零；若$f$为常数，结论就是假设本身。设$f=\sum_{i=0}^m a_ix^i$，$a_m\ne0$、$m\ge1$，且$f^N=0$。$f^N$的$x^{mN}$系数为$a_m^N$，故$a_m^N=0$。

令$I=(a_m)$。因环交换，每个$I$中元素都为$ra_m$，任意$N$个这类元素相乘均为零，所以$I^N=0$。在商环$(R/I)[x]$中，$f$的像$\bar f$仍满足$\bar f^{\,N}=0$，而最高项已经消失。若$\bar f=0$，每个系数都属于$I$；否则其次数至多为$m-1$，由归纳假设，其每个系数均幂零。因此两种情形中，对每个$i<m$都存在$k_i\ge1$使
\[
a_i^{k_i}\in I.
\]
由于$I^N=0$，有$(a_i^{k_i})^N=a_i^{k_iN}=0$，故$a_i$幂零。连同$a_m$的幂零性，归纳完成。

\end{proof}

\paragraph{复习提示}

商环中幂零只给出$a_i^{k_i}\in I$，须利用$I^N=0$才能提升回母环的幂零性。



\subsection{习题44：多项式零因子的常数湮灭元判据}

\label{ex:3:44}

\paragraph{题目}

设$R$为交换幺环，$0\ne f\in R[x]$。证明$f$为零因子当且仅当存在$0\ne a\in R$使$af=0$。这里补明非零条件，以符合原笔记零因子的定义。

\paragraph{考点与思路}

沿原解取最低次的非零湮灭多项式，逐一消去$f$的最高系数，再取一个非零系数作常数湮灭元。

\paragraph{所用结论与原文依据}

\noteref{zerodiv}；\noteref{c3poly}；\noteref{ringlaws}。原解位置：\path{chapters/Part 2/chapter3-exercises.tex}，第470行起。本题在现有习题解答的思路内补齐推导。

% 原题和原解来源同上；新增中间推导为本册补写。

\begin{proof}

本题要求$f\ne0$，因为本笔记的零因子按定义非零。充分性为：若$0\ne a\in R$且$af=0$，将$a$视为非零常数多项式，即得到$f$与非零元素的零乘积，故$f$为零因子。

再证必要性。设$f=\sum_{i=0}^m a_ix^i\ne0$，$a_m\ne0$。因$f$为零因子，存在非零多项式$g$使$fg=0$。在全部这类$g$中选次数$d$最小者，写$g=\sum_{j=0}^d b_jx^j$，$b_d\ne0$。

先由$fg=0$的$x^{m+d}$系数得$a_mb_d=0$。因此$a_mg$若非零，其次数小于$d$；而交换性给出$f(a_mg)=a_m(fg)=0$，与$d$最小矛盾。所以$a_mg=0$，从而$(f-a_mx^m)g=0$。

现作向下归纳。若已知$a_{k+1}g=\cdots=a_mg=0$，则
\[
\left(\sum_{i=0}^k a_ix^i\right)g=0.
\]
比较$x^{k+d}$系数得$a_kb_d=0$，故$a_kg$若非零则次数小于$d$。同时$f(a_kg)=a_k(fg)=0$，仍违背最小性。因此$a_kg=0$。由归纳，对每个$i$都有$a_ig=0$。

比较这些等式的$x^d$系数，得$a_ib_d=0$对全部$i$成立。取非零常数$a=b_d$，便有$af=0$。

\end{proof}

\paragraph{复习提示}

原题若包括$f=0$，右侧在非零环中成立而左侧不成立，所以必须补明$f\ne0$。



\subsection{习题45：常数项可逆且其余系数幂零的多项式可逆}

\label{ex:3:45}

\paragraph{题目}

设$R$为交换环，$f=a_0x^n+a_1x^{n-1}+\cdots+a_n$。若$a_n$是单位而$a_0,\ldots,a_{n-1}$均幂零，证明$f$是$R[x]$的单位。

\paragraph{考点与思路}

沿原解除去可逆常数项，把剩余多项式证明为幂零，再用有限等比和给出逆。

\paragraph{所用结论与原文依据}

\noteref{unit}；\noteref{c3nil}；\noteref{ringlaws}。原解位置：\path{chapters/Part 2/chapter3-exercises.tex}，第477行起。本题在现有习题解答的思路内补齐推导。

% 原题和原解来源同上；新增中间推导为本册补写。

\begin{proof}

设$f=a_0x^n+\cdots+a_{n-1}x+a_n$，$a_n$为单位。令
\[
h=a_n^{-1}(f-a_n)=\sum_{i=0}^{n-1}(a_n^{-1}a_i)x^{n-i}.
\]
每个$a_i$幂零，取$m_i\ge1$使$a_i^{m_i}=0$。交换性给出$(a_n^{-1}a_i)^{m_i}=0$。若$n=0$，$f=a_n$已经是单位；以下设$n\ge1$，并令$N=1+\sum_{i=0}^{n-1}(m_i-1)$。

将$h^N$按分配律展开。每项选取$N$个系数因子；若第$i$种系数均只出现至多$m_i-1$次，总数将小于$N$。所以每项至少含某个$(a_n^{-1}a_i)^{m_i}$，从而该项为零。因此$h^N=0$。

令$S=\sum_{j=0}^{N-1}(-h)^j$。有限和相乘消去中间项，得到
\[
(1+h)S=S(1+h)=1-(-h)^N=1.
\]
由于$f=a_n(1+h)$，故
\[
\boxed{f^{-1}=a_n^{-1}\sum_{j=0}^{N-1}(-h)^j.}
\]

\end{proof}

\paragraph{复习提示}

有限多个幂零系数保证存在一个共同的展开次数；不能只说“每个系数幂零，所以和幂零”而不说明理由。



\subsection{习题46：各系数幂零推出多项式幂零}

\label{ex:3:46}

\paragraph{题目}

证明习题43的逆命题。

\paragraph{考点与思路}

沿原解给出统一指数，并明确展开中使用的抽屉计数。

\paragraph{所用结论与原文依据}

\noteref{c3nil}；\exref{3}{43}。原解位置：\path{chapters/Part 2/chapter3-exercises.tex}，第488行起。本题在现有习题解答的思路内补齐推导。

% 原题和原解来源同上；新增中间推导为本册补写。

\begin{proof}

设$f=\sum_{i=0}^m a_ix^i$，并对每个$i$选$n_i\ge1$使$a_i^{n_i}=0$。令
\[
N=1+\sum_{i=0}^m(n_i-1).
\]
将$f^N$按分配律展开，每项都是$N$个$a_ix^i$的乘积。设第$i$类因子出现$k_i$次，则$\sum_i k_i=N$。若所有$k_i<n_i$，则$\sum_i k_i\le\sum_i(n_i-1)=N-1$，矛盾。因此至少有一个$i$满足$k_i\ge n_i$。

环交换，故该项的系数可重排为$\prod_i a_i^{k_i}$，其中含某个为零的因子$a_i^{n_i}$。所以每项为零，$f^N=0$。连同\exref{3}{43}，得到$f$幂零当且仅当其每个系数幂零。

\end{proof}

\paragraph{复习提示}

指数$N$只需保证充分大，不必是$f$的最小幂零指数。



\subsection{习题47：多元多项式环的幂零根}

\label{ex:3:47}

\paragraph{题目}

设$R$为交换幺环，$R'=R[x_1,\ldots,x_n]$，$r(R)$为$R$的幂零根。证明$r(R')=r(R)[x_1,\ldots,x_n]$。补明本组题目沿用的交换性条件。

\paragraph{考点与思路}

按原解逐个添加变量，反复使用一元系数判据。

\paragraph{所用结论与原文依据}

\noteref{c3nil}；\noteref{c3multi}；\exref{3}{43}；\exref{3}{46}。原解位置：\path{chapters/Part 2/chapter3-exercises.tex}，第495行起。本题在现有习题解答的思路内补齐推导。

% 原题和原解来源同上；新增中间推导为本册补写。

\begin{proof}

这里$R$为交换幺环，$r(R)$表示幂零元全体，也等于全部素理想之交。由\exref{3}{43}与\exref{3}{46}，对$f\in R[x]$，有
\[
f\in r(R[x])\ \Longleftrightarrow\ f\text{的每个系数属于 }r(R).
\]
所以$r(R[x])=r(R)[x]$，右侧指所有系数均属于$r(R)$的多项式集合。

对变量个数$n$归纳。$n=1$已证，$n=0$时也是恒等式。若对$n-1$成立，令$S=R[x_1,\ldots,x_{n-1}]$。它仍是交换幺环，故
\begin{align*}
r(R[x_1,\ldots,x_n])
&=r(S[x_n])\\
&=r(S)[x_n]\\
&=\bigl(r(R)[x_1,\ldots,x_{n-1}]\bigr)[x_n]\\
&=r(R)[x_1,\ldots,x_n].
\end{align*}
最后一步是把关于$x_n$的各个系数再次按前$n-1$个变量展开；最终每个标量系数都在$r(R)$中，反向也成立。

\end{proof}

\paragraph{复习提示}

右侧是系数属于幂零根的多项式理想，不是宣称整个理想有统一幂零指数。



\subsection{习题48：有限环中右逆必为左逆}

\label{ex:3:48}

\paragraph{题目}

证明有限环中右逆也是左逆，并判断无限环是否仍成立。

\paragraph{考点与思路}

保留有限集上乘法映射的证明，并把无限维空间上的移位算子逐个定义。

\paragraph{所用结论与原文依据}

\noteref{unit}；\exref{3}{2}；\exref{3}{3}。原解位置：\path{chapters/Part 2/chapter3-exercises.tex}，第502行起。本题在现有习题解答的思路内补齐推导。

% 原题和原解来源同上；新增中间推导为本册补写。

\begin{proof}

设幺环$R$有限，且$uv=1$。对每个$y\in R$，有$u(vy)=y$，所以$L_u:x\mapsto ux$满射。有限集合上的满射必单射。又
\[
u(vu)=(uv)u=u=u1,
\]
单射性给出$vu=1$。因此右逆也是左逆。

无限环中不成立。取域$F$上以$e_0,e_1,e_2,\ldots$为基的向量空间$V$（每个向量为有限线性组合），令$R=\operatorname{End}_F(V)$，乘法取映射复合。定义线性映射
\[
v(e_i)=e_{i+1}\ (i\ge0),\qquad u(e_0)=0,\quad u(e_{i+1})=e_i\ (i\ge0).
\]
给定基上的像唯一决定线性映射。对每个$i\ge0$，$(uv)(e_i)=e_i$，所以$uv=\mathrm{id}_V$。但$(vu)(e_0)=0\ne e_0$，故$vu\ne\mathrm{id}_V$。$R$是无限幺环，这给出所需反例。

\end{proof}

\paragraph{复习提示}

无限维向量空间用有限支撑的向量即可定义移位，不需要引入无限和。



\subsection{习题49：多项式环中强拟正则元就是幂零元}

\label{ex:3:49}

\paragraph{题目}

证明交换环$R$的多项式环$R[x]$中，强拟正则元恰为幂零元。

\paragraph{考点与思路}

保留原解从$1-xf$可逆得到系数幂零的路线；补证原解省略的“单位多项式正次数系数幂零”。

\paragraph{所用结论与原文依据}

\noteref{c3nil}；\noteref{c3prime}；\noteref{c3degree}；\exref{3}{22}；\exref{3}{45}；\exref{3}{46}。原解位置：\path{chapters/Part 2/chapter3-exercises.tex}，第509行起。本题在现有习题解答的思路内补齐推导。

% 原题和原解来源同上；新增中间推导为本册补写。

\begin{proof}

先证所需的一般事实：若$u=\sum a_ix^i$是交换幺环$R$上多项式环的单位，则每个$i>0$的$a_i$幂零。对任一素理想$P$，逐系数取模把$u$送为整环$(R/P)[x]$中的单位。整环上非零乘积次数可加，所以该单位只能是常数，因此$a_i\in P$对所有$i>0$成立。

因$P$任意，这些系数属于全部素理想之交$r(R)$，即均幂零。若$R$为零环，系数全为零，此结论也成立。此外，单位的常数项必为$R$的单位，因为把逆元等式取常数项即可。连同\exref{3}{45}，这也补齐了多项式单位的双向系数判据。

现设$f$在$R[x]$中强拟正则。定义要求对每个$g\in R[x]$，$1-fg$可逆，特别取$g=x$，则$1-xf$可逆。若$f=\sum a_ix^i$，则$1-xf$的$x^{i+1}$系数为$-a_i$，由刚证事实，每个$a_i$都幂零。根据\exref{3}{46}，$f$幂零。

反过来，若$f^N=0$，对任意$g\in R[x]$，交换性给出$(fg)^N=f^Ng^N=0$。于是
\[
(1-fg)\sum_{j=0}^{N-1}(fg)^j
=\sum_{j=0}^{N-1}(fg)^j(1-fg)=1-(fg)^N=1.
\]
所以$1-fg$对每个$g$均为单位，即$f$强拟正则。

\end{proof}

\paragraph{复习提示}

单位多项式的必要系数条件不能只引用第45题；第45题只证明了充分性，此处已补证必要性。



\subsection{习题50：每个元素满足某个幂等关系时素理想极大}

\label{ex:3:50}

\paragraph{题目}

设交换环$R$中每个$a$都有某个$n>1$使$a^n=a$。证明每个素理想均极大。

\paragraph{考点与思路}

在素理想商整环内使用幂关系，并以无零因子性获得逆元。

\paragraph{所用结论与原文依据}

\noteref{c3prime}；\noteref{field}。原解位置：\path{chapters/Part 2/chapter3-exercises.tex}，第522行起。本题在现有习题解答的思路内补齐推导。

% 原题和原解来源同上；新增中间推导为本册补写。

\begin{proof}

设$P$为$R$的素理想，$D=R/P$为非零整环。任取$0\ne\bar a\in D$，选代表元$a\in R$。题设给出某个$n>1$使$a^n=a$，故
\[
\bar a^n=\bar a,\qquad \bar a(\bar a^{n-1}-1)=0.
\]
因$\bar a\ne0$且$D$无零因子，得$\bar a^{n-1}=1$。于是$\bar a^{n-2}$是$\bar a$的逆元；当$n=2$时按零次幂为$1$理解。
因此$D$每个非零元均可逆，是域，故$P$为极大理想。指数$n$允许随$a$变化，证明只对当前元素选取其对应的$n$。

\end{proof}

\paragraph{复习提示}

不能在原环中直接消去$a$，因为原环可以有零因子；应先到素理想的商整环。



\subsection{习题51：布尔环的特征、素商与有限生成理想}

\label{ex:3:51}

\paragraph{题目}

设$R$为布尔环，即每个$a$满足$a^2=a$。证明$2a=0$、每个素理想极大且商域为$\mathbb F_2$、每个有限生成理想均为主理想。

\paragraph{考点与思路}

沿原解由$(a+a)^2$得到$2a=0$，再展开$(a+b)^2$得到交换性，最后合并两个理想生成元。

\paragraph{所用结论与原文依据}

\noteref{c3prime}；\noteref{c3principal}；\noteref{ringlaws}。原解位置：\path{chapters/Part 2/chapter3-exercises.tex}，第529行起。本题在现有习题解答的思路内补齐推导。

% 原题和原解来源同上；新增中间推导为本册补写。

\begin{proof}

\begin{enumerate}
\item 对任意$a\in R$，布尔条件给出$a^2=a$及$(a+a)^2=a+a$。分配律展开后有$4a^2=2a$，即$4a=2a$，在加法群中相减得$2a=0$。特别地，$-a=a$。

对$a+b$再用布尔条件，得$a^2+ab+ba+b^2=a+b$，故$ab+ba=0$。由$-ba=ba$，得到$ab=ba$，所以$R$交换。
\item 在本组素理想题目的约定下，$R$有幺元。设$P$为素理想，则$D=R/P$是非零整环，且每个元素$z$仍满足$z^2=z$。由$z(z-1)=0$，每个$z$只能为$0$或$1$。二者因商环非零而不同，所以$D$恰有两个元素，其加法、乘法分别由$1+1=0$和$1\cdot1=1$确定。因此$D\cong\F_2$，且$P$极大。
\item 任取$a,b\in R$，令$c=a+b+ab$。交换性、$a^2=a,b^2=b$及$2ab=0$给出
\[
ac=a^2+ab+a^2b=a,\qquad bc=ab+b^2+ab^2=b.
\]
所以$a,b\in(c)$，得到$(a,b)\subseteq(c)$。另一方面$c\in(a,b)$，故$(c)\subseteq(a,b)$，于是$(a,b)=(c)$。

对生成元个数归纳：一个生成元的理想已是主理想；若$(a_1,\ldots,a_k)=(d)$，则$(a_1,\ldots,a_{k+1})=(d,a_{k+1})$可按上述公式再合并成一个生成元。零理想为$(0)$，所以全部有限生成理想均为主理想。
\end{enumerate}

\end{proof}

\paragraph{复习提示}

先得到$2a=0$再推出交换性，不能在尚未证明交换时把$(a+b)^2$写成$a^2+2ab+b^2$。



\subsection{习题52：唯一素理想、非单位幂零与根商域的等价性}

\label{ex:3:52}

\paragraph{题目}

设$R\ne0$为交换幺环，$r(R)$为其幂零根。证明以下三项等价：(1) $R$只有一个素理想；(2) $R$的每个元素不是单位就是幂零元；(3) $R/r(R)$是域。这里补明非零条件。

\paragraph{考点与思路}

保留原解的三个条件循环推导，并在最后一步用商环的素理想对应补齐唯一性。

\paragraph{所用结论与原文依据}

\noteref{c3nil}；\noteref{c3prime}；\noteref{c3maxexists}；\exref{3}{20}。原解位置：\path{chapters/Part 2/chapter3-exercises.tex}，第540行起。本题在现有习题解答的思路内补齐推导。

% 原题和原解来源同上；新增中间推导为本册补写。

\begin{proof}

这里须设$R\ne0$为交换幺环。将三条件分别记为(1)、(2)、(3)。若允许零环，(2)成立而(1)、(3)均不成立，等价性会失败。

\begin{enumerate}
\item[(1)$\Rightarrow$(2)] 设唯一素理想为$P$。由幂零根等于全部素理想之交，有$r(R)=P$。$R\ne0$保证至少有一个极大理想，而每个极大理想均为素理想，所以$P$本身是极大理想，且是唯一的极大理想。

若$a\notin P$但不可逆，由\exref{3}{20}，$a$属于某个极大理想，只能是$P$，矛盾。因此$P$外的元素都是单位；$P=r(R)$内的元素都是幂零元。得到(2)。
\item[(2)$\Rightarrow$(3)] 非零环的单位不可能幂零：若$a^m=0$且$a$可逆，乘以$a^{-m}$会得$1=0$。所以$r(R)$不含$1$，为真理想。取$R/r(R)$中任意非零类$a+r(R)$，则$a\notin r(R)$，不是幂零元。由(2)，$a$为单位，故$a^{-1}+r(R)$为该类的逆。于是$R/r(R)$为域。
\item[(3)$\Rightarrow$(1)] 商环为域使$r(R)$为极大理想，特别也是素理想。对任意素理想$P$，由$r(R)$为全部素理想之交，有$r(R)\subseteq P$。因为$r(R)$极大且$P\ne R$，只能$P=r(R)$。因此素理想唯一。
\end{enumerate}

\end{proof}

\paragraph{复习提示}

最后一步需要先由商域得到$r(R)$极大，而不能从“每个素理想包含根”直接推出相等。



\subsection{习题53：只有有限多个理想的整环是域}

\label{ex:3:53}

\paragraph{题目}

证明只有有限多个理想的整环是域。

\paragraph{考点与思路}

保留主理想幂链的路线，先从有限性推出相邻相等，再用整环消去。

\paragraph{所用结论与原文依据}

\noteref{c3principal}；\noteref{zerodiv}；\noteref{field}。原解位置：\path{chapters/Part 2/chapter3-exercises.tex}，第549行起。本题在现有习题解答的思路内补齐推导。

% 原题和原解来源同上；新增中间推导为本册补写。

\begin{proof}

任取$0\ne a\in R$。考虑下降链
\[
(a)\supseteq(a^2)\supseteq(a^3)\supseteq\cdots.
\]
因$R$只有有限多个理想，这些理想不能全部不同，所以存在$i<j$使$(a^i)=(a^j)$。下降关系给出
\[
(a^i)\supseteq(a^{i+1})\supseteq(a^j)=(a^i),
\]
故$(a^i)=(a^{i+1})$。于是$a^i\in(a^{i+1})$，在交换幺环中可写$a^i=ba^{i+1}$，即
\[
a^i(1-ba)=0.
\]
整环中$a\ne0$使$a^i\ne0$，故$1-ba=0$。交换性给出$ab=ba=1$。因此任一非零元$a$可逆，$R$为域。

\end{proof}

\paragraph{复习提示}

有限多个理想不直接意味着任选某一级就相邻相等，需要利用重复及链的包含关系。



\subsection{习题54：理想集合有极小元的整环是域}

\label{ex:3:54}

\paragraph{题目}

设整环$R$的任意非空理想集合按包含关系都有极小元。证明$R$是域。

\paragraph{考点与思路}

沿原解在指定的幂主理想集合中取极小元，再展开与上一题相同的消去。

\paragraph{所用结论与原文依据}

\noteref{c3principal}；\exref{3}{53}。原解位置：\path{chapters/Part 2/chapter3-exercises.tex}，第556行起。本题在现有习题解答的思路内补齐推导。

% 原题和原解来源同上；新增中间推导为本册补写。

\begin{proof}

固定$0\ne a\in R$。集合$\mathcal A=\{(a^m):m\ge1\}$非空，由题设存在其中一个极小元$(a^m)$。因$(a^{m+1})\in\mathcal A$且$(a^{m+1})\subseteq(a^m)$，极小性给出相等。

因此有$b\in R$使$a^m=ba^{m+1}$，即$a^m(1-ba)=0$。$R$为整环且$a\ne0$，所以$a^m\ne0$，得到$ba=1$；交换性使$ab=1$。故每个非零元均可逆，$R$为域。

\end{proof}

\paragraph{复习提示}

“极小元”只须在这一个理想集合内使用，不需要把它当成全部非零理想中的最小理想。



\subsection{习题55：华罗庚恒等式中的可逆性与因子次序}

\label{ex:3:55}

\paragraph{题目}

设$a,b,ab-1$均为单位。证明$a-b^{-1}$和$(a-b^{-1})^{-1}-a^{-1}$均为单位，并验证华罗庚恒等式。

\paragraph{考点与思路}

严格保留原解的分解$c=(ab-1)b^{-1}$，以及$c^{-1}-a^{-1}$的三因子分解。

\paragraph{所用结论与原文依据}

\noteref{unit}；\noteref{ringlaws}。原解位置：\path{chapters/Part 2/chapter3-exercises.tex}，第563行起。本题在现有习题解答的思路内补齐推导。

% 原题和原解来源同上；新增中间推导为本册补写。

\begin{proof}

令$c=a-b^{-1}$。因$a,b,ab-1$均可逆，
\[
c=(ab-1)b^{-1}
\]
是两个单位之积，所以可逆；具体逆元为$bc_1$，其中$c_1=(ab-1)^{-1}$。

接着保持顺序计算
\begin{align*}
c^{-1}(a-c)a^{-1}
&=c^{-1}aa^{-1}-c^{-1}ca^{-1}\\
&=c^{-1}-a^{-1}.
\end{align*}
由于$a-c=b^{-1}$，得到
\[
c^{-1}-a^{-1}=c^{-1}b^{-1}a^{-1}.
\]
右侧为三个单位的乘积，故也是单位。按反序取逆，
\begin{align*}
\bigl((a-b^{-1})^{-1}-a^{-1}\bigr)^{-1}
&=(c^{-1}b^{-1}a^{-1})^{-1}\\
&=abc\\
&=ab(a-b^{-1})=aba-a.
\end{align*}
这就是所要求的华罗庚恒等式，亦可写为$aba=a+((a-b^{-1})^{-1}-a^{-1})^{-1}$。每次取逆的对象都已先证明可逆，计算没有交换因子。

\end{proof}

\paragraph{复习提示}

非交换环中取乘积的逆要颠倒顺序；$abc$不能改成任意其他排列。



\subsection{习题56：共轭不变除子环的中心性}

\label{ex:3:56}

\paragraph{题目}

设$D$为除环，$C$为其中心，$R$为子体且对每个$0\ne d\in D$有$dRd^{-1}\subseteq R$。证明$R=D$或$R\subseteq C$。

\paragraph{考点与思路}

将第12题的$x,x+1$比较法推广到任意除环，保留原解顺序。

\paragraph{所用结论与原文依据}

\noteref{field}；\noteref{c3anti}；\exref{3}{12}。原解位置：\path{chapters/Part 2/chapter3-exercises.tex}，第577行起。本题在现有习题解答的思路内补齐推导。

% 原题和原解来源同上；新增中间推导为本册补写。

\begin{proof}

记$C=Z(D)=\{z\in D:zd=dz\text{ 对每个 }d\in D\}$。若$R=D$，结论成立。以下设$R\ne D$，取$x\in D\setminus R$。除子环$R$含$0,1$，所以$x,x+1$都非零。

任取$a\in R$，题设给出$b=xax^{-1}\in R$和$c=(x+1)a(x+1)^{-1}\in R$。有$xa=bx$及$xa+a=cx+c$，故
\[
(b-c)x=c-a.
\]
若$b\ne c$，$R$中存在$(b-c)^{-1}$，从而$x=(b-c)^{-1}(c-a)\in R$，矛盾。故$b=c$，进而$c=a$，于是$xa=ax$。该论证对每个$x\notin R$都成立。

再取$y\in R$，固定一个$x\notin R$。因$x+y\notin R$，它也与$a$交换，所以
\[
ya=(x+y)a-xa=a(x+y)-ax=ay.
\]
因此$a$与$R$内外的每个元素都交换，即$a\in C$。$a$任意，得$R\subseteq C$。

\end{proof}

\paragraph{复习提示}

除子环对非零元素的取逆封闭，是由$b-c\ne0$推出$x\in R$的关键。



\subsection{习题57：实多项式函数环到实数的同态}

\label{ex:3:57}

\paragraph{题目}

设$R$为实系数多项式函数环。证明每个$t\in\R$的代入映射$\eta_t(f)=f(t)$为到$\R$的满同态，且每个保幺同态$R\to\R$（因而每个满同态）均唯一为某个$\eta_t$。一般不保幺同态中的零映射需另列。

\paragraph{考点与思路}

沿原解先用正数为平方证明常数上的保序性，再用有理数稠密性固定所有实数。

\paragraph{所用结论与原文依据}

\noteref{c3functions1}；\noteref{c3poly}。原解位置：\path{chapters/Part 2/chapter3-exercises.tex}，第586行起。本题在现有习题解答的思路内补齐推导。

% 原题和原解来源同上；新增中间推导为本册补写。

\begin{proof}

将自变量函数记为$s:\R\to\R$，$s(c)=c$。原笔记说明无限域上的不同多项式给出不同多项式函数，因此$R=\R[s]\cong\R[x]$。

对$t\in\R$，映射$\eta_t(f)=f(t)$满足
\[
\eta_t(f+g)=f(t)+g(t),\quad
\eta_t(fg)=f(t)g(t),\quad\eta_t(1)=1.
\]
每个实数$c$是常函数$c$的像，所以$\eta_t$满射。

反过来，设$\eta:R\to\R$为保幺同态（特别地，任意满同态都保幺）。限制在常数域上，若$c\ne0$，由$\eta(c)\eta(c^{-1})=1$知$\eta(c)\ne0$。若$c>0$，写$c=d^2$且$d\ne0$，则$\eta(c)=\eta(d)^2>0$。因此对$c_1<c_2$，$\eta(c_2)-\eta(c_1)=\eta(c_2-c_1)>0$，故该限制严格保序。

保幺和加法性使$\eta$固定所有整数，保逆使它固定所有有理数。任给$c\in\R$，若$\eta(c)>c$，由有理数稠密性选$c<r<\eta(c)$，$r\in\Q$；保序性却给出$\eta(c)<\eta(r)=r$，矛盾。若$\eta(c)<c$，选$\eta(c)<r<c$，则$\eta(r)<\eta(c)$也矛盾。故$\eta(c)=c$对所有常数成立。

令$t=\eta(s)$，则对$f=\sum_i a_is^i$，
\[
\eta(f)=\sum_i\eta(a_i)\eta(s)^i=\sum_i a_it^i=f(t).
\]
所以$\eta=\eta_t$，且$t$由$s$的像唯一确定。若不要求保幺或满射，零同态也是环同态，并不属于上述代入映射；因此原题中“这样的同态”须作此区分。

\end{proof}

\paragraph{复习提示}

不能预设实数常数都被固定；须先由平方、保序和有理数稠密性证明。



\subsection{习题58：无限域上非零多项式的非零点集}

\label{ex:3:58}

\paragraph{题目}

设$F$为无限域。若$g\ne0$，且$f$在所有满足$g(c_1,\ldots,c_n)\ne0$的点上为零，证明$f=0$。

\paragraph{考点与思路}

保留原解先看乘积$fg$的路线，并展开无限域上多元零函数只能来自零多项式的归纳。

\paragraph{所用结论与原文依据}

\noteref{c3functionsN}；\noteref{c3degree}；\noteref{c3roots}。原解位置：\path{chapters/Part 2/chapter3-exercises.tex}，第596行起。本题在现有习题解答的思路内补齐推导。

% 原题和原解来源同上；新增中间推导为本册补写。

\begin{proof}

设$f,g\in F[x_1,\ldots,x_n]$，$g\ne0$。任取$a\in F^n$。若$g(a)\ne0$，题设给出$f(a)=0$；若$g(a)=0$，乘积值也为零。因此
\[
(fg)(a)=0\qquad(a\in F^n).
\]
下面沿原笔记按变量个数归纳，说明无限域上的多项式若处处取零值，则为零多项式。$n=1$时，非零多项式的根数至多为次数，不可能有无限多个根。

设结论对$n-1$成立，写$h=\sum_{j=0}^m h_j(x_1,\ldots,x_{n-1})x_n^j$，并假设$h$处处为零。固定前$n-1$个坐标，得到关于$x_n$的一元多项式，在无限个点上为零，故各系数值$h_j(a_1,\ldots,a_{n-1})$全为零。因前$n-1$个坐标任意，由归纳假设各$h_j=0$，所以$h=0$。

将该结论用于$h=fg$，得$fg=0$作为多项式恒等式。域上的多元多项式环是整环（逐次将一个变量加到整环上，非零首项系数的乘积保持非零）。因$g\ne0$，得到$f=0$。

\end{proof}

\paragraph{复习提示}

“处处取零值”和“零多项式”的等价性依赖无限域，须先证明再用于乘积。



\subsection{习题59：有限域上原点的指示多项式}

\label{ex:3:59}

\paragraph{题目}

设$F$有$q$个元素。证明$g_0=\prod_{i=1}^n(1-x_i^{q-1})$在原点取值$1$，在其余点取值$0$。

\paragraph{考点与思路}

沿原解逐个因子判断，并交代$q-1$次幂公式来自有限乘法群。

\paragraph{所用结论与原文依据}

\noteref{lagrange}；\noteref{c3finite}。原解位置：\path{chapters/Part 2/chapter3-exercises.tex}，第603行起。本题在现有习题解答的思路内补齐推导。

% 原题和原解来源同上；新增中间推导为本册补写。

\begin{proof}

有限域$F$的非零元素构成$q-1$阶乘法群，所以由拉格朗日定理，对每个$a\ne0$都有$a^{q-1}=1$。又$q\ge2$，故$q-1>0$，$0^{q-1}=0$。

在原点，$g_0$的每个因子$1-x_i^{q-1}$均取$1$，所以
\[
g_0(0,\ldots,0)=1.
\]
在任意非零点$(a_1,\ldots,a_n)$，至少有一个$a_i\ne0$，则其对应因子$1-a_i^{q-1}=0$，整个乘积为零。

因此$g_0$恰在原点取$1$，在$F^n\setminus\{0\}$上取$0$。

\end{proof}

\paragraph{复习提示}

这里$n\ge1$，且有限域至少两个元素，因而所有出现的$q-1$次幂均为正次幂。



\subsection{习题60：只有原点为零点的多项式次数下界}

\label{ex:3:60}

\paragraph{题目}

若$f$在原点为零而在其余点均非零，证明$g=1-f^{q-1}$与习题59的$g_0$取值相同，并推出$\deg f\geqslant n$。

\paragraph{考点与思路}

按原解先化为原点指示函数，再逐变量降次，并展开低于$q$次数代表的唯一性。

\paragraph{所用结论与原文依据}

\noteref{c3functionsN}；\noteref{c3finite}；\exref{3}{59}。原解位置：\path{chapters/Part 2/chapter3-exercises.tex}，第610行起。本题在现有习题解答的思路内补齐推导。

% 原题和原解来源同上；新增中间推导为本册补写。

\begin{proof}

设$f\in F[x_1,\ldots,x_n]$，$f(0)=0$而$f(a)\ne0$对每个$a\ne0$成立。令$g=1-f^{q-1}$。由有限域非零元的$q-1$次幂为$1$，$g$在原点取$1$，其余点取$0$；故与\exref{3}{59}的$g_0$给出同一函数。

对任意变量$x_i$及$e\ge q$，有
\[
x_i^e-x_i^{e-q+1}=x_i^{e-q}(x_i^q-x_i).
\]
右侧在$F^n$上恒为零，因此可在不改变函数的前提下把指数$e$降低$q-1$。逐个单项式、逐个变量反复降低，得到代表元$r$，使每个变量的次数均小于$q$。每次降低不增加单项式的总次数，合并同类项也不增加总次数，所以$\deg r\le\deg g$。

还须说明这种代表唯一。若$r$对每个变量次数均小于$q$且处处为零，对变量个数归纳。一个变量时，次数小于$q$的非零多项式不能有$q$个根。多个变量时，写$r=\sum_{j=0}^{q-1}r_j(x_1,\ldots,x_{n-1})x_n^j$。固定前$n-1$个坐标后，用一元结论使每个$r_j$的值均为零，再用归纳假设得$r_j=0$。因此$r=0$，两个给出同一函数的低次数代表必相等。

将$g$降次所得$r$与$g_0$比较，因$g_0$自身对每个变量的次数小于$q$，唯一性给出$r=g_0$。而$g_0$的单项式$x_1^{q-1}\cdots x_n^{q-1}$系数为$(-1)^n\ne0$，故$\deg g_0=n(q-1)$。于是
\[
n(q-1)=\deg r\le\deg g\le(q-1)\deg f.
\]
因$q-1>0$，得到$\deg f\ge n$。这里次数均为总次数。

\end{proof}

\paragraph{复习提示}

同一函数不保证多项式相等；只有逐变量次数小于$q$后，才能使用代表唯一性。



\subsection{习题61：低于变量数次数的多项式还有非零零点}

\label{ex:3:61}

\paragraph{题目}

设$F$是$q$元域，$\deg f<n$且$f(0)=0$。证明$f$还有一个非零零点。

\paragraph{考点与思路}

保留原解的有限域求和计数，补证幂和消失以及整数零点数与域元素的区别。

\paragraph{所用结论与原文依据}

\noteref{c3finite}；\noteref{c3functionsN}；\exref{3}{59}。原解位置：\path{chapters/Part 2/chapter3-exercises.tex}，第617行起。本题在现有习题解答的思路内补齐推导。

% 原题和原解来源同上；新增中间推导为本册补写。

\begin{proof}

若$f=0$，任取非零向量即为所求零点。以下设$f\ne0$，记$N$为$f$在$F^n$中的零点个数，$p=\operatorname{char}F$。在域$F$中求和，有
\[
N\cdot1_F=\sum_{a\in F^n}\bigl(1-f(a)^{q-1}\bigr),
\]
因为一个求和项在零点取$1$，非零点取$0$。

先计算$S_k=\sum_{t\in F}t^k$，其中$k=0$时求和项按常数函数$1$理解，故$S_0=q\cdot1_F=0$。对$k>0$，取有限乘法群$F^\times$的生成元$\zeta$。若$q-1\nmid k$，则$\zeta^k\ne1$；乘以$\zeta$置换非零元素，给出$S_k=\zeta^k S_k$，故$S_k=0$。若$q-1\mid k$，则每个非零$t$都满足$t^k=1$，所以$S_k=(q-1)\cdot1_F=-1_F$。

对单项式$x_1^{k_1}\cdots x_n^{k_n}$，逐坐标求和的分配律给出
\[
\sum_{a\in F^n}a_1^{k_1}\cdots a_n^{k_n}=\prod_{i=1}^n S_{k_i}.
\]
此乘积可能非零，必须每个$k_i$都是$q-1$的正倍数，特别是总次数至少为$n(q-1)$。但$f^{q-1}$的总次数为$(q-1)\deg f<n(q-1)$，所以其每个单项式在全体点上的和均为零。常数$1$的总和也是$q^n\cdot1_F=0$。故$N\cdot1_F=0$，即整数$N$被$p$整除。

原点已经是零点，所以$N\ge1$。正整数$N$被$p$整除，故$N\ge p\ge2$，必有另一个不同于原点的零点。

\end{proof}

\paragraph{复习提示}

域中得到的等式是$N\cdot1_F=0$，不是整数$N=0$；它说明零点个数被特征整除。



\subsection{习题62：形式幂级数的逆、整环性与商域表示}

\label{ex:3:62}

\paragraph{题目}

设$R$为交换环、$f=\sum_{i\geqslant0}a_ix^i\in R[[x]]$。证明：(i) $f$可逆当且仅当$a_0$可逆；(ii) $R[[x]]$为整环当且仅当$R$为整环；(iii) 若$R$为整环且$F$为其商域，则$R[[x]]$商域的非零元均可写作$x^\nu(a_0+a_1x+\cdots)$，其中$\nu\in\mathbb Z$、$a_i\in F$、$a_0\ne0$。

\paragraph{考点与思路}

沿原解递推逆元系数，再以最低非零项检验整环性，最后直接展开任意两个幂级数之商。

\paragraph{所用结论与原文依据}

\noteref{c3fraction}；\noteref{c3poly}；\noteref{unit}。原解位置：\path{chapters/Part 2/chapter3-exercises.tex}，第628行起。本题在现有习题解答的思路内补齐推导。

% 原题和原解来源同上；新增中间推导为本册补写。

\begin{proof}

先说明形式幂级数的运算。$R[[x]]$由全部序列$(a_0,a_1,\ldots)$组成，写成$\sum_{i\ge0}a_ix^i$；相等是逐系数相等，加法逐系数进行，乘法的$x^m$系数定义为$\sum_{i=0}^m a_ib_{m-i}$。每个系数仅用有限次运算，不涉及解析收敛。交换幺环$R$的运算律逐系数保证$R[[x]]$仍为交换幺环。

\begin{enumerate}
\item 若$fg=1$，比较常数项得$a_0b_0=1$，所以$a_0$可逆。反过来，若$a_0$为单位，递推定义
\[
b_0=a_0^{-1},\qquad
b_m=-a_0^{-1}\sum_{i=1}^m a_i b_{m-i}\quad(m\ge1).
\]
当定义$b_m$时，右侧只涉及已定义的$b_0,\ldots,b_{m-1}$，故此序列存在且唯一。令$g=\sum b_mx^m$，则$fg$的常数项为$1$，其余第$m$项系数为
\[
a_0b_m+\sum_{i=1}^m a_i b_{m-i}=0.
\]
因此$fg=1$；交换性使$gf=1$，故$f$可逆。
\item 若$R$为整环，设非零$f,g$的最低非零项分别为$a_rx^r,b_sx^s$。$fg$中低于$r+s$次的系数全零，第$r+s$次系数恰为$a_rb_s\ne0$，因此$fg\ne0$。常数$1$非零，所以$R[[x]]$为整环。

反过来，常数嵌入$R\hookrightarrow R[[x]]$为单射，保加法、乘法和单位元。若$R[[x]]$是整环，则其常数子环非零且无零因子，故$R$为整环。
\item 设$F$为$R$的商域，逐系数嵌入$R[[x]]\hookrightarrow F[[x]]$。任取非零分式$f/g$，其中$f,g\in R[[x]]\setminus\{0\}$。分离最低非零项，写成$f=x^r u$、$g=x^s v$，其中$u,v\in R[[x]]$的常数项非零。在$F[[x]]$中，这两个常数项都是单位，故由第一项，$u,v$均可逆。于是
\[
\frac fg=x^{r-s}uv^{-1}=x^\nu(a_0+a_1x+a_2x^2+\cdots),
\]
其中$\nu=r-s\in\Z$，$a_i\in F$，且$a_0=u(0)/v(0)\ne0$。

这可在形式洛朗级数域$F((x))$中理解；它由仅有有限个负次幂项的级数组成。非零洛朗级数分离最低次幂后，其余因子常数项非零，可按第一项求逆，所以确实构成域。上述分式嵌入良定义且单射：$f/g=f'/g'$等价于$fg'=f'g$，而逐系数嵌入保此等式；其像为零只可能$f=0$。

因此证明的是$R[[x]]$商域的每个非零元都有上述表示，并非任意$F$系数洛朗级数都必来自该商域。当$R=F$时，任一洛朗级数均可写为$x^{-m}h$，$h\in F[[x]]$，故$\operatorname{Frac}(F[[x]])=F((x))$。
\end{enumerate}

\end{proof}

\paragraph{复习提示}

形式级数不需要收敛；一般$R$时只得到其商域嵌入$F((x))$，不能无条件宣称二者相等。



\subsection{习题63：非单位构成唯一极大理想的条件}

\label{ex:3:63}

\paragraph{题目}

设$R\ne0$为交换幺环，且每个非单位$a$都使$1+a$为单位。证明全部非单位构成唯一的极大理想，并判断这种环是否存在。补明非零条件。

\paragraph{考点与思路}

沿原解把假设的单位和归一化为$1$，用$1+a$的条件排除两个非单位之和为单位。

\paragraph{所用结论与原文依据}

\noteref{c3prime}；\exref{3}{30}。原解位置：\path{chapters/Part 2/chapter3-exercises.tex}，第645行起。本题在现有习题解答的思路内补齐推导。

% 原题和原解来源同上；新增中间推导为本册补写。

\begin{proof}

补明$R\ne0$为交换幺环。令$M=R\setminus R^\times$，则$0\in M$而$1\notin M$。若$a\in M$、$r\in R$且$ra$是单位，设$(ra)c=1$，交换性给出$a(rc)=1$，使$a$为单位，矛盾。因此$ra\in M$。若$-a$可逆，则$a=(-1)(-a)$也可逆，故$-a\in M$。

任取$a,b\in M$。若$a+b=u$为单位，则$au^{-1},bu^{-1}$均不是单位；否则再乘单位$u$会使$a$或$b$成为单位。于是$-au^{-1}$为非单位，题设给出
\[
1-au^{-1}=bu^{-1}\in R^\times,
\]
矛盾。因此$a+b\in M$。结合上述性质，$M$是理想，且因不含$1$而为真理想。

任一真理想都不能含单位，所以均包含于$M$。若$M\subsetneq J\lhd R$，则$J$含一个不属于$M$的单位，故$J=R$。所以$M$极大；任一极大理想既含于$M$，就只能等于$M$，故唯一。

这种环存在。例如$R=\Z/(p^n)$，$n\ge1$：非单位恰为$p$的倍数，若$p\mid a$，则$p\nmid1+a$，所以$1+\bar a$为单位。唯一极大理想是$(p)/(p^n)$。

\end{proof}

\paragraph{复习提示}

归一化时仍须确认$au^{-1}$非单位；零环需要排除，否则非单位集为空而非理想。



\subsection{习题64：单位加零构成子域及其外元素超越性}

\label{ex:3:64}

\paragraph{题目}

设交换整环$R$中每个非零非单位$a$都使$1+a$仍非单位。证明全部单位加上$0$构成子域$F$，且$R\setminus F$中的元素都超越于$F$。

\paragraph{考点与思路}

沿原解用两个单位之和归一化制造矛盾，再把代数关系化为常数项非零的关系。

\paragraph{所用结论与原文依据}

\noteref{field}；\noteref{c3algebraic}；\noteref{unit}。原解位置：\path{chapters/Part 2/chapter3-exercises.tex}，第652行起。本题在现有习题解答的思路内补齐推导。

% 原题和原解来源同上；新增中间推导为本册补写。

\begin{proof}

令$F=R^\times\cup\{0\}$。整环要求$1\ne0$，故$F$含$0,1$。单位的积及逆仍为单位，单位的负元也是单位，所以$F$对乘法、取负和非零元取逆均封闭。

只需检查加法。若一个加数为零，无需再证。设$u,v$均为单位，且$u+v\ne0$。反设$u+v$不是单位，令$a=-(u+v)u^{-1}$。乘以单位保持非单位性，且整环中乘单位保持非零，故$a$是非零非单位。题设要求$1+a$也不是单位，但
\[
1+a=1-(u+v)u^{-1}=-vu^{-1}
\]
是单位，矛盾。因此$u+v$不是零就是单位，$F$对加法封闭，是$R$的子域。

任取$z\in R\setminus F$，则$z\ne0$且不是单位。若$z$在$F$上代数，存在非零多项式$\sum_{i=0}^m c_iX^i$使$\sum c_iz^i=0$。取最小的$c_k\ne0$，因$R$为整环且$z\ne0$，消去公共因子$z^k$，得到常数项非零的关系。重新编号可写为
\[
c_0+c_1z+\cdots+c_dz^d=0,\qquad c_0\ne0.
\]
此处$d\ge1$，否则非零常数为零。因$c_0$是$F$中的非零元，故在$R$中可逆，且
\[
z\bigl[-c_0^{-1}(c_1+c_2z+\cdots+c_dz^{d-1})\bigr]=1.
\]
交换性使括号内元素为$z$的双侧逆，矛盾。因此$z$在$F$上超越。

\end{proof}

\paragraph{复习提示}

若最初关系常数项为零，要利用$z\ne0$和整环消去，不能直接假定有非零常数项。



\subsection{习题65：无限域上的乘法型加法恒等式无非常数多项式解}

\label{ex:3:65}

\paragraph{题目}

设$F$为无限域。证明不存在非常数多项式$f\in F[x]$满足$f(x)f(y)=f(x+y)$。

\paragraph{考点与思路}

沿原解比较二元多项式的最高总次数；若题意是逐点等式，先用无限域保证它是恒等式。

\paragraph{所用结论与原文依据}

\noteref{c3degree}；\noteref{c3functionsN}；\exref{3}{58}。原解位置：\path{chapters/Part 2/chapter3-exercises.tex}，第665行起。本题在现有习题解答的思路内补齐推导。

% 原题和原解来源同上；新增中间推导为本册补写。

\begin{proof}

反设$f\in F[x]$非常数，次数为$d>0$，首项系数为$c\ne0$。若题目给的是对每个$x,y\in F$的等式，则多项式
\[
H(X,Y)=f(X)f(Y)-f(X+Y)
\]
在$F^2$上处处为零。因$F$无限，由原笔记的多元多项式函数定理（其归纳证明也在\exref{3}{58}展开），$H$为零多项式。若题目本来就指形式恒等式，这一步直接成立。

但$f(X)f(Y)$的单项式$X^dY^d$系数为$c^2\ne0$，总次数为$2d$。另一方面，$f(X+Y)$的每项由$(X+Y)^j$、$j\le d$展开而来，总次数至多为$d$，所以其中$X^dY^d$的系数为零。因$2d>d$，两边系数不可能相等，矛盾。

\end{proof}

\paragraph{复习提示}

无限域条件用于把逐点函数等式提升为多项式恒等式；形式恒等式的次数比较本身不需要无限性。



\subsection{习题66：全矩阵环的左理想都是主左理想}

\label{ex:3:66}

\paragraph{题目}

设$F$为域，证明$M_n(F)$的每个左理想都是主理想。

\paragraph{考点与思路}

保留原解的行空间构造，证明任意出现过的行都能单独放进任何行位置，再以一组基组成生成矩阵。

\paragraph{所用结论与原文依据}

\noteref{ideal}；\noteref{subring}。原解位置：\path{chapters/Part 2/chapter3-exercises.tex}，第672行起。本题在现有习题解答的思路内补齐推导。

% 原题和原解来源同上；新增中间推导为本册补写。

\begin{proof}

左理想$L\subseteq M_n(F)$是加法子群，且对任意$X\in M_n(F),A\in L$有$XA\in L$。设$W$为$L$中所有矩阵的所有行向量组成的集合。下面先证它为$F^n$的子空间，而不预先假定这些行都来自同一个矩阵。

若$w$为$A\in L$的第$i$行，则$E_{ki}A\in L$，其第$k$行为$w$、其他行为零。于是任意两个出现过的行$w,z$都能放到同一个指定行后相加；乘以标量矩阵$cE_n$还能把行变为$cw$。零矩阵给出零行，因此$W$为子空间。

记$L_W$为每行都属于$W$的矩阵集合。按定义$L\subseteq L_W$。反过来，若$B\in L_W$，对其每行$w_k$，选一个$L$中矩阵使其某行等于$w_k$，再左乘矩阵单位把它搬到第$k$行、其余行置零。将这$n$个矩阵相加，得到$B\in L$。故$L=L_W$。

取$W$的一组基$w_1,\ldots,w_r$，$0\le r\le n$，令$A$的前$r$行为这些基向量，其余行为零。因$A\in L_W=L$，有$M_n(F)A\subseteq L$。任给$B\in L$，每行$b_i$可写为$\sum_{j=1}^r c_{ij}w_j$。令$X$的前$r$列为这些$c_{ij}$，其余列为零，则$XA=B$。因此$L\subseteq M_n(F)A$，得到
\[
L=M_n(F)A.
\]
当$W=0$时取$A=0$，同样成立。故每个左理想都是主左理想。

\end{proof}

\paragraph{复习提示}

左乘矩阵形成原矩阵各行的线性组合，因此左理想对应行空间；不能在此换成列空间。



\subsection{习题67：阿贝尔群的自同态环}

\label{ex:3:67}

\paragraph{题目}

设$G$为阿贝尔群，证明$\operatorname{End}(G)$以逐点加法和映射复合构成含幺环。

\paragraph{考点与思路}

按原解逐项验证两种运算，特别写出依赖自同态加法性的分配律。

\paragraph{所用结论与原文依据}

\noteref{ring}；\noteref{hom}；\noteref{aut}。原解位置：\path{chapters/Part 2/chapter3-exercises.tex}，第679行起。本题在现有习题解答的思路内补齐推导。

% 原题和原解来源同上；新增中间推导为本册补写。

\begin{proof}

将$G$的群运算写为加法。对$\eta,\xi\in\operatorname{End}(G)$定义
\[
(\eta+\xi)(a)=\eta(a)+\xi(a),\qquad
(\eta\xi)(a)=\eta(\xi(a)).
\]
若$a,b\in G$，因$G$交换，
\begin{align*}
(\eta+\xi)(a+b)
&=\eta(a)+\eta(b)+\xi(a)+\xi(b)\\
&=(\eta(a)+\xi(a))+(\eta(b)+\xi(b)).
\end{align*}
所以$\eta+\xi$仍为自同态。又$(\eta\xi)(a+b)=\eta(\xi(a)+\xi(b))=\eta\xi(a)+\eta\xi(b)$，故乘法也封闭。

逐点加法由$G$的交换加法继承结合律、交换律；零映射为零元，$(-\eta)(a)=-\eta(a)$为加法逆元，且$-\eta$仍为自同态。因此加法构成阿贝尔群。映射复合本来结合，恒等映射$\mathrm{id}_G$为双侧乘法单位。

对任意$a\in G$，有
\begin{align*}
[\eta(\xi+\zeta)](a)
&=\eta(\xi(a)+\zeta(a))=\eta\xi(a)+\eta\zeta(a),\\
[(\eta+\xi)\zeta](a)
&=\eta(\zeta(a))+\xi(\zeta(a))=\eta\zeta(a)+\xi\zeta(a).
\end{align*}
映射相等可逐点判断，故左右分配律成立，$\operatorname{End}(G)$为幺环。当$G=0$时，它是允许$0=1$的零环。

\end{proof}

\paragraph{复习提示}

阿贝尔条件用于保证逐点和仍为群同态；乘法是复合，不是逐点相乘。



\subsection{习题68：循环群、有理数群与初等阿贝尔群的自同态环}

\label{ex:3:68}

\paragraph{题目}

分别求$\mathbb Z^+$、无限循环群、$\mathbb Q^+$及$n$个$p$阶循环群的直和的自同态环。

\paragraph{考点与思路}

依原解顺序，用生成元的像、$1$的像以及一组$p$元基的像确定全部自同态，并检查环运算。

\paragraph{所用结论与原文依据}

\noteref{cyclic}；\noteref{hom}；\exref{3}{67}。原解位置：\path{chapters/Part 2/chapter3-exercises.tex}，第691行起。本题在现有习题解答的思路内补齐推导。

% 原题和原解来源同上；新增中间推导为本册补写。

\begin{proof}

各环的乘法均按\exref{3}{67}取复合，右侧映射先作用。
\begin{enumerate}
\item 对整数加法群$\Z^+$，任一自同态$\eta$由$m=\eta(1)\in\Z$唯一决定，因为对所有$k\in\Z$，$\eta(k)=km$。反过来，$\eta_m(k)=mk$是自同态。于是$m\mapsto\eta_m$为双射，且
\[
\eta_m+\eta_n=\eta_{m+n},\qquad
\eta_m\eta_n=\eta_{mn},\qquad \eta_1=\mathrm{id}.
\]
故$\operatorname{End}(\Z^+)\cong\Z$。
\item 若$G=\langle g\rangle$为无限循环群，任一自同态由$\eta(g)=g^m$、$m\in\Z$决定，即$\eta_m(g^k)=g^{mk}$。因群交换，写成加法记号后，两个自同态相加对应指数相加，复合对应指数相乘。因此同样有$\operatorname{End}(G)\cong\Z$。这里“自同态相加”使用群运算；若用乘法记号，则$(\eta+\xi)(x)=\eta(x)\xi(x)$。
\item 对有理数加法群$\Q^+$，设$r=\eta(1)\in\Q$。任给$q=m/n$，$m\in\Z,n>0$，加法性使
\[
n\eta(m/n)=\eta(m)=mr.
\]
在$\Q$中除以$n$，得到$\eta(q)=rq$，故由$r$唯一决定。反过来，每个乘法映射$\mu_r(q)=rq$都是加法自同态。它们满足$\mu_r+\mu_s=\mu_{r+s}$、$\mu_r\mu_s=\mu_{rs}$及$\mu_1=\mathrm{id}$。因此$\operatorname{End}(\Q^+)\cong\Q$。
\item 设$G=\bigoplus\limits_{i=1}^n C_p$，各因子选一个生成元$e_i$，改用加法记号。每个元素唯一写成$\sum_i a_i e_i$，$a_i\in\F_p$。以$\bar m\,g=mg$定义标量乘法，因$pg=0$，不依赖整数代表元，所以$G$为$\F_p$上的$n$维向量空间。

每个群自同态$T$满足$T(mg)=mT(g)$，因此自动为$\F_p$线性映射；每个线性映射反过来也是群自同态。写$T(e_j)=\sum_i t_{ij}e_i$，以列向量记录坐标，得到矩阵$A_T=(t_{ij})$。每个矩阵都唯一决定这样的线性映射，故$T\mapsto A_T$为双射。并且
\[
A_{T+S}=A_T+A_S,\qquad A_{TS}=A_TA_S,\qquad
A_{\mathrm{id}}=E_n.
\]
因此$\operatorname{End}(G)\cong M_n(\F_p)$，且此处是通常矩阵乘法的同构。若$n=0$，两边均理解为零环。
\end{enumerate}

\end{proof}

\paragraph{复习提示}

取列向量约定时，右侧映射先作用恰对应$A_{TS}=A_TA_S$；不要误写成反同构。




\clearpage\section*{原笔记引用定位}\addcontentsline{toc}{section}{原笔记引用定位}
下表只记录原笔记条目的位置，不重抄定理和定义。原笔记未启用前两章时，可直接按源文件行号回查。正文引用链接指向本表；每项均保留原章节的实际编号或节内位置。
\small\begin{longtable}{p{6.3cm}p{8.0cm}}
\toprule 原笔记条目 & 原文位置与内容范围\\\midrule\endhead

\phantomsection\label{note:group}原笔记第1章第2节，定义 1.2.1 & 群的定义；第1章“群的定义和简单性质”。\newline\path{chapters/Part 2/chapter1.tex}，第51行起。\\

\phantomsection\label{note:glaws}原笔记第1章第2节，命题 1.2.3 & 单位元、逆元与反序取逆；第1章“群的定义和简单性质”。\newline\path{chapters/Part 2/chapter1.tex}，第77行起。\\

\phantomsection\label{note:cancel}原笔记第1章第2节，命题 1.2.4 & 方程、消去律；第1章“群的定义和简单性质”。\newline\path{chapters/Part 2/chapter1.tex}，第90行起。\\

\phantomsection\label{note:powers}原笔记第1章第2节，命题 1.2.5 & 整数指数法则；第1章“群的定义和简单性质”。\newline\path{chapters/Part 2/chapter1.tex}，第109行起。\\

\phantomsection\label{note:abelianpowers}原笔记第1章第2节，命题 1.2.7 & 交换元素的乘积幂；第1章“群的定义和简单性质”。\newline\path{chapters/Part 2/chapter1.tex}，第134行起。\\

\phantomsection\label{note:dihedral}原笔记第1章第3节，定义 1.3.2 & 二面体群及其阶数与乘法关系；第1章“群的例子”。\newline\path{chapters/Part 2/chapter1.tex}，第163行起。\\

\phantomsection\label{note:cycles}原笔记第1章第3节，定义 1.3.7 & 轮换与对换；第1章“群的例子”。\newline\path{chapters/Part 2/chapter1.tex}，第208行起。\\

\phantomsection\label{note:cycledecomp}原笔记第1章第3节，定理 1.3.9 & 不相交轮换分解；第1章“群的例子”。\newline\path{chapters/Part 2/chapter1.tex}，第222行起。\\

\phantomsection\label{note:subgroup}原笔记第1章第4节，定理 1.4.3 & 子群判别法；第1章“子群、陪集”。\newline\path{chapters/Part 2/chapter1.tex}，第260行起。\\

\phantomsection\label{note:cosets}原笔记第1章第4节，命题 1.4.8 & 陪集相等的判据与陪集划分；第1章“子群、陪集”。\newline\path{chapters/Part 2/chapter1.tex}，第315行起。\\

\phantomsection\label{note:lagrange}原笔记第1章第4节，定理 1.4.10 & 拉格朗日定理；第1章“子群、陪集”。\newline\path{chapters/Part 2/chapter1.tex}，第342行起。\\

\phantomsection\label{note:order}原笔记第1章第4节，命题 1.4.12 & 元素的阶与幂的判据；第1章“子群、陪集”。\newline\path{chapters/Part 2/chapter1.tex}，第357行起。\\

\phantomsection\label{note:iso}原笔记第1章第5节，定义 1.5.1 & 群同构及同构不变量；第1章“群的同构”。\newline\path{chapters/Part 2/chapter1.tex}，第387行起。\\

\phantomsection\label{note:cayley}原笔记第1章第5节，定理 1.5.3 & 凯莱定理及左正则表示；第1章“群的同构”。\newline\path{chapters/Part 2/chapter1.tex}，第406行起。\\

\phantomsection\label{note:sign}原笔记第1章第4节，命题 1.4.6 & 符号同态与交错群的阶；第1章“子群、陪集”。\newline\path{chapters/Part 2/chapter1.tex}，第293行起。\\

\phantomsection\label{note:hom}原笔记第1章第6节，定义 1.6.1 & 群同态、像与核；第1章“同态、正规子群”。\newline\path{chapters/Part 2/chapter1.tex}，第481行起。\\

\phantomsection\label{note:normal}原笔记第1章第6节，命题 1.6.5 & 正规子群的等价判据；第1章“同态、正规子群”。\newline\path{chapters/Part 2/chapter1.tex}，第577行起。\\

\phantomsection\label{note:setproduct}原笔记第1章第7节，命题 1.7.2 & 集合乘积及取逆公式；第1章“商群”。\newline\path{chapters/Part 2/chapter1.tex}，第615行起。\\

\phantomsection\label{note:quotient}原笔记第1章第7节，定义 1.7.4 & 商群及自然同态；第1章“商群”。\newline\path{chapters/Part 2/chapter1.tex}，第680行起。\\

\phantomsection\label{note:firstiso}原笔记第1章第7节，定理 1.7.7 & 群同态基本定理；第1章“商群”。\newline\path{chapters/Part 2/chapter1.tex}，第730行起。\\

\phantomsection\label{note:ring}原笔记第1章第8节，定义 1.8.1 & 环的定义；第1章“环、子环”。\newline\path{chapters/Part 2/chapter1.tex}，第780行起。\\

\phantomsection\label{note:ringlaws}原笔记第1章第8节，命题 1.8.2 & 环中零元、负元和分配律的运算；第1章“环、子环”。\newline\path{chapters/Part 2/chapter1.tex}，第790行起。\\

\phantomsection\label{note:subring}原笔记第1章第8节，定理 1.8.5 & 子环判别法；第1章“环、子环”。\newline\path{chapters/Part 2/chapter1.tex}，第889行起。\\

\phantomsection\label{note:ringiso}原笔记第1章第8节，定义 1.8.7 & 环同构及其性质；第1章“环、子环”。\newline\path{chapters/Part 2/chapter1.tex}，第922行起。\\

\phantomsection\label{note:unit}原笔记第1章第9节，定义 1.9.1 & 幺环、单位及双侧逆元；第1章“各种特殊类型的环”。\newline\path{chapters/Part 2/chapter1.tex}，第953行起。\\

\phantomsection\label{note:zerodiv}原笔记第1章第9节，定义 1.9.3 & 左右零因子、无零因子的消去律；第1章“各种特殊类型的环”。\newline\path{chapters/Part 2/chapter1.tex}，第987行起。\\

\phantomsection\label{note:field}原笔记第1章第9节，定义 1.9.6 & 域与体；第1章“各种特殊类型的环”。\newline\path{chapters/Part 2/chapter1.tex}，第1033行起。\\

\phantomsection\label{note:quaternion}原笔记第1章第9节，命题 1.9.9 & 四元数体的矩阵模型；第1章“各种特殊类型的环”。\newline\path{chapters/Part 2/chapter1.tex}，第1091行起。\\

\phantomsection\label{note:ideal}原笔记第1章第10节，定义 1.10.3 & 理想的定义；第1章“环的同态、理想”。\newline\path{chapters/Part 2/chapter1.tex}，第1200行起。\\

\phantomsection\label{note:idealunit}原笔记第1章第10节，命题 1.10.5 & 含有单位的理想；第1章“环的同态、理想”。\newline\path{chapters/Part 2/chapter1.tex}，第1234行起。\\

\phantomsection\label{note:correspondence}原笔记第2章第1节，定理 2.1.1 & 子群对应定理；第2章“群的同态定理”。\newline\path{chapters/Part 2/chapter2.tex}，第18行起。\\

\phantomsection\label{note:secondiso}原笔记第2章第1节，定理 2.1.3 & 群的第二同构定理；第2章“群的同态定理”。\newline\path{chapters/Part 2/chapter2.tex}，第111行起。\\

\phantomsection\label{note:generated}原笔记第2章第1节 & 生成子群与有限生成群；第2章“群的同态定理”。\newline\path{chapters/Part 2/chapter2.tex}，第142行起。\\

\phantomsection\label{note:derived}原笔记第2章第1节，定义 2.1.6 & 换位子与换位子群；第2章“群的同态定理”。\newline\path{chapters/Part 2/chapter2.tex}，第172行起。\\

\phantomsection\label{note:abelianization}原笔记第2章第1节，定理 2.1.7 & 交换同态像与换位子群；第2章“群的同态定理”。\newline\path{chapters/Part 2/chapter2.tex}，第176行起。\\

\phantomsection\label{note:cyclic}原笔记第2章第2节，定理 2.2.2 & 循环群及其子群分类；第2章“循环群”。\newline\path{chapters/Part 2/chapter2.tex}，第228行起。\\

\phantomsection\label{note:powerorder}原笔记第2章第2节，命题 2.2.3 & 幂的阶的公式；第2章“循环群”。\newline\path{chapters/Part 2/chapter2.tex}，第254行起。\\

\phantomsection\label{note:coprime}原笔记第2章第2节，引理 2.2.4 & 互素阶交换元素乘积的阶；第2章“循环群”。\newline\path{chapters/Part 2/chapter2.tex}，第276行起。\\

\phantomsection\label{note:simple}原笔记第2章第3节，定义 2.3.1 & 单群的定义；第2章“单群与交错群的单性”。\newline\path{chapters/Part 2/chapter2.tex}，第331行起。\\

\phantomsection\label{note:threecycles}原笔记第2章第3节，引理 2.3.3 & 对换与3-轮换的生成结论；第2章“单群与交错群的单性”。\newline\path{chapters/Part 2/chapter2.tex}，第347行起。\\

\phantomsection\label{note:conjugation}原笔记第2章第3节，命题 2.3.4 & 轮换的共轭公式；第2章“单群与交错群的单性”。\newline\path{chapters/Part 2/chapter2.tex}，第366行起。\\

\phantomsection\label{note:an}原笔记第2章第3节，定理 2.3.6 & 交错群的单性；第2章“单群与交错群的单性”。\newline\path{chapters/Part 2/chapter2.tex}，第398行起。\\

\phantomsection\label{note:derivedseries}原笔记第2章第4节，定义 2.4.1 & 导出列；第2章“可解群”。\newline\path{chapters/Part 2/chapter2.tex}，第503行起。\\

\phantomsection\label{note:solvabledef}原笔记第2章第4节，定义 2.4.2 & 可解群的定义；第2章“可解群”。\newline\path{chapters/Part 2/chapter2.tex}，第521行起。\\

\phantomsection\label{note:solvable}原笔记第2章第4节，命题 2.4.8 & 可解群的子群、像与扩张；第2章“可解群”。\newline\path{chapters/Part 2/chapter2.tex}，第681行起。\\

\phantomsection\label{note:aut}原笔记第2章第5节，定义 2.5.1 & 自同构与自同构群；第2章“群的自同构群”。\newline\path{chapters/Part 2/chapter2.tex}，第717行起。\\

\phantomsection\label{note:cyclicaut}原笔记第2章第5节，命题 2.5.3 & 循环群的自同构；第2章“群的自同构群”。\newline\path{chapters/Part 2/chapter2.tex}，第740行起。\\

\phantomsection\label{note:inner}原笔记第2章第5节，定理 2.5.6 & 内自同构群与中心商群；第2章“群的自同构群”。\newline\path{chapters/Part 2/chapter2.tex}，第819行起。\\

\phantomsection\label{note:center}原笔记第2章第5节，定义 2.5.5 & 群的中心；第2章“群的自同构群”。\newline\path{chapters/Part 2/chapter2.tex}，第807行起。\\

\phantomsection\label{note:autcenterless}原笔记第2章第5节，定理 2.5.8 & 无中心群的自同构定理；第2章“群的自同构群”。\newline\path{chapters/Part 2/chapter2.tex}，第854行起。\\

\phantomsection\label{note:complete}原笔记第2章第5节，定义 2.5.9 & 完全群；第2章“群的自同构群”。\newline\path{chapters/Part 2/chapter2.tex}，第884行起。\\

\phantomsection\label{note:characteristic}原笔记第2章第5节，定义 2.5.10 & 特征子群及其性质；第2章“群的自同构群”。\newline\path{chapters/Part 2/chapter2.tex}，第890行起。\\

\phantomsection\label{note:action}原笔记第2章第6节，定义 2.6.1 & 群作用及对应的置换表示；第2章“群在集合上的作用”。\newline\path{chapters/Part 2/chapter2.tex}，第923行起。\\

\phantomsection\label{note:orbit}原笔记第2章第6节，定理 2.6.11 & 轨道稳定子定理；第2章“群在集合上的作用”。\newline\path{chapters/Part 2/chapter2.tex}，第1060行起。\\

\phantomsection\label{note:pcenter}原笔记第2章第6节，推论 2.6.16 & 有限$p$群的中心非平凡；第2章“群在集合上的作用”。\newline\path{chapters/Part 2/chapter2.tex}，第1140行起。\\

\phantomsection\label{note:centerizer}原笔记第2章第6节，定义 2.6.17 & 元素的中心化子；第2章“群在集合上的作用”。\newline\path{chapters/Part 2/chapter2.tex}，第1161行起。\\

\phantomsection\label{note:classeq}原笔记第2章第6节，定理 2.6.18 & 类方程；第2章“群在集合上的作用”。\newline\path{chapters/Part 2/chapter2.tex}，第1165行起。\\

\phantomsection\label{note:normalizer}原笔记第2章第6节，定义 2.6.20 & 子群的正规化子；第2章“群在集合上的作用”。\newline\path{chapters/Part 2/chapter2.tex}，第1207行起。\\

\phantomsection\label{note:sylow1}原笔记第2章第7节，定理 2.7.2 & 西罗第一定理；第2章“西罗定理”。\newline\path{chapters/Part 2/chapter2.tex}，第1259行起。\\

\phantomsection\label{note:sylowdef}原笔记第2章第7节，定义 2.7.3 & 西罗子群的定义；第2章“西罗定理”。\newline\path{chapters/Part 2/chapter2.tex}，第1286行起。\\

\phantomsection\label{note:sylow2}原笔记第2章第7节，定理 2.7.4 & 西罗第二定理；第2章“西罗定理”。\newline\path{chapters/Part 2/chapter2.tex}，第1290行起。\\

\phantomsection\label{note:sylow3}原笔记第2章第7节，定理 2.7.7 & 西罗第三定理；第2章“西罗定理”。\newline\path{chapters/Part 2/chapter2.tex}，第1341行起。\\

\phantomsection\label{note:internal}原笔记第2章第8节，定理 2.8.3 & 内直和分解定理；第2章“群的直和”。\newline\path{chapters/Part 2/chapter2.tex}，第1542行起。\\

\phantomsection\label{note:free}原笔记第2章第11节，定理 2.11.7 & 自由群的泛性质与生成关系；第2章“自由幺半群与自由群”。\newline\path{chapters/Part 2/chapter2.tex}，第1958行起。\\

\phantomsection\label{note:c3correspondence}原笔记第3章“环的同态定理”节 & 环的对应定理；现为源码注释，按源码回查。\newline\path{chapters/Part 2/chapter3.tex}，第130行起。\\
\phantomsection\label{note:c3isomorphism}原笔记第3章“环的同态定理”节 & 第三同构定理及其证明；现为源码注释，按源码回查。\newline\path{chapters/Part 2/chapter3.tex}，第169行起。\\
\phantomsection\label{note:c3direct}原笔记第3章“环的直和”节 & 直和定义、分量理想与内直和；现为源码注释，按源码回查。\newline\path{chapters/Part 2/chapter3.tex}，第254行起。\\
\phantomsection\label{note:c3product}原笔记第3章“环的直和”节 & 理想积的定义与分配律；现为源码注释，按源码回查。\newline\path{chapters/Part 2/chapter3.tex}，第332行起。\\
\phantomsection\label{note:c3comaximal}原笔记第3章“环的直和”节 & 互素理想的积等于交；现为源码注释，按源码回查。\newline\path{chapters/Part 2/chapter3.tex}，第356行起。\\
\phantomsection\label{note:c3crt}原笔记第3章“环的直和”节 & 中国剩余定理及证明；现为源码注释，按源码回查。\newline\path{chapters/Part 2/chapter3.tex}，第379行起。\\
\phantomsection\label{note:c3principal}原笔记第3章“环的直和”节 & 生成理想与主理想的定义及表达式；现为源码注释，按源码回查。\newline\path{chapters/Part 2/chapter3.tex}，第442行起。\\
\phantomsection\label{note:c3anti}原笔记第3章“环的反同构”节 & 反同构、反同态与反自同构定义；现为源码注释，按源码回查。\newline\path{chapters/Part 2/chapter3.tex}，第469行起。\\
\phantomsection\label{note:c3conjugate}原笔记第3章“环的反同构”节 & 四元数共轭与乘法表；现为源码注释，按源码回查。\newline\path{chapters/Part 2/chapter3.tex}，第498行起。\\
\phantomsection\label{note:c3prime}原笔记第3章“素理想和极大理想”节 & 素理想、极大理想及商环判据；现为源码注释，按源码回查。\newline\path{chapters/Part 2/chapter3.tex}，第564行起。\\
\phantomsection\label{note:c3maxexists}原笔记第3章“素理想和极大理想”节 & 极大理想存在性与佐恩引理构造；现为源码注释，按源码回查。\newline\path{chapters/Part 2/chapter3.tex}，第708行起。\\
\phantomsection\label{note:c3nil}原笔记第3章“素理想和极大理想”节 & 幂零元等于全部素理想之交；现为源码注释，按源码回查。\newline\path{chapters/Part 2/chapter3.tex}，第717行起。\\
\phantomsection\label{note:c3fraction}原笔记第3章“商域和分式环”节 & 商域的存在与分式表示；当前正文。\newline\path{chapters/Part 2/chapter3.tex}，第758行起。\\
\phantomsection\label{note:c3poly}原笔记第3章“交换环上的多项式环”节 & 多项式环的定义、系数唯一性和代入；当前正文。\newline\path{chapters/Part 2/chapter3.tex}，第1149行起。\\
\phantomsection\label{note:c3multi}原笔记第3章“交换环上的多项式环”节 & 多元多项式环的定义；当前正文。\newline\path{chapters/Part 2/chapter3.tex}，第1296行起。\\
\phantomsection\label{note:c3algebraic}原笔记第3章“交换环上的多项式环”节 & 代数关系与代数无关的定义；当前正文。\newline\path{chapters/Part 2/chapter3.tex}，第1340行起。\\
\phantomsection\label{note:c3polyquot}原笔记第3章“交换环上的多项式环”节 & 系数取商诱导的多项式商环同构；当前正文。\newline\path{chapters/Part 2/chapter3.tex}，第1366行起。\\
\phantomsection\label{note:c3degree}原笔记第3章“整环上的一元多项式环”节 & 次数法则、整环性与单位；现为源码注释，按源码回查。\newline\path{chapters/Part 2/chapter3.tex}，第1392行起。\\
\phantomsection\label{note:c3division}原笔记第3章“整环上的一元多项式环”节 & 带余除法与因式定理；现为源码注释，按源码回查。\newline\path{chapters/Part 2/chapter3.tex}，第1429行起。\\
\phantomsection\label{note:c3roots}原笔记第3章“整环上的一元多项式环”节 & 非零多项式的根数上界；现为源码注释，按源码回查。\newline\path{chapters/Part 2/chapter3.tex}，第1473行起。\\
\phantomsection\label{note:c3qroots}原笔记第3章“整环上的一元多项式环”节 & 四元数二次方程的共轭根构造；现为源码注释，按源码回查。\newline\path{chapters/Part 2/chapter3.tex}，第1490行起。\\
\phantomsection\label{note:c3finite}原笔记第3章“整环上的一元多项式环”节 & 有限乘法子群循环、有限域非零元的幂；现为源码注释，按源码回查。\newline\path{chapters/Part 2/chapter3.tex}，第1517行起。\\
\phantomsection\label{note:c3minimal}原笔记第3章“整环上的一元多项式环”节 & 极小多项式定义与商环刻画；现为源码注释，按源码回查。\newline\path{chapters/Part 2/chapter3.tex}，第1569行起。\\
\phantomsection\label{note:c3irred}原笔记第3章“整环上的一元多项式环”节 & 不可约定义与极大理想、商域的等价性；现为源码注释，按源码回查。\newline\path{chapters/Part 2/chapter3.tex}，第1587行起。\\
\phantomsection\label{note:c3nilideal}原笔记第3章“整环上的一元多项式环”节 & 幂零理想定义；现为源码注释，按源码回查。\newline\path{chapters/Part 2/chapter3.tex}，第1699行起。\\
\phantomsection\label{note:c3functions1}原笔记第3章“多项式函数”节 & 一元多项式函数与形式多项式的对应；现为源码注释，按源码回查。\newline\path{chapters/Part 2/chapter3.tex}，第1746行起。\\
\phantomsection\label{note:c3functionsN}原笔记第3章“多项式函数”节 & 无限域的多元零函数判据及有限域低次数代表的唯一性；现为源码注释，按源码回查。\newline\path{chapters/Part 2/chapter3.tex}，第1802行起。\\
\bottomrule\end{longtable}\normalsize

\clearpage\section*{题号与证明入口}\addcontentsline{toc}{section}{题号与证明入口}
\small\begin{longtable}{cp{13.5cm}}
\toprule 章 & 题号及主题（点击跳转）\\\midrule\endhead

1 & \hyperref[ex:1:1]{1：平方公式判别交换群}\\

1 & \hyperref[ex:1:2]{2：每个元素的平方为单位元}\\

1 & \hyperref[ex:1:3]{3：有限消去半群是群}\\

1 & \hyperref[ex:1:4]{4：可解方程的结合系统}\\

1 & \hyperref[ex:1:5]{5：非交换平方公式的反例}\\

1 & \hyperref[ex:1:6]{6：任意偶数阶的非交换例子}\\

1 & \hyperref[ex:1:7]{7：反复共轭}\\

1 & \hyperref[ex:1:8]{8：取逆映射的同态性}\\

1 & \hyperref[ex:1:9]{9：由子集定义等价关系}\\

1 & \hyperref[ex:1:10]{10：$n\Z\cong\Z$}\\

1 & \hyperref[ex:1:11]{11：克莱因四元群}\\

1 & \hyperref[ex:1:12]{12：指数2的子群正规}\\

1 & \hyperref[ex:1:13]{13：偶数阶群有2阶元素}\\

1 & \hyperref[ex:1:14]{14：二面体群的矩阵实现}\\

1 & \hyperref[ex:1:15]{15：连续三个指数的乘积公式}\\

1 & \hyperref[ex:1:16]{16：有限阶元素的积可以无限阶}\\

1 & \hyperref[ex:1:17]{17：交换群的有限阶元素成子群}\\

1 & \hyperref[ex:1:18]{18：只有有限多个子群的群必有限}\\

1 & \hyperref[ex:1:19]{19：$D_n$的全部正规子群}\\

1 & \hyperref[ex:1:20]{20：子群乘积的判别}\\

1 & \hyperref[ex:1:21]{21：有限子群乘积的计数}\\

1 & \hyperref[ex:1:22]{22：正规子群的乘积}\\

1 & \hyperref[ex:1:23]{23：中心化子与正规化子}\\

1 & \hyperref[ex:1:24]{24：2阶群}\\

1 & \hyperref[ex:1:25]{25：4阶群的分类}\\

1 & \hyperref[ex:1:26]{26：素数阶群}\\

1 & \hyperref[ex:1:27]{27：整数对上的环}\\

1 & \hyperref[ex:1:28]{28：交换整数加法与乘法的角色}\\

1 & \hyperref[ex:1:29]{29：改变坐标得到同构环}\\

1 & \hyperref[ex:1:30]{30：子环与母环的单位元和交换性}\\

1 & \hyperref[ex:1:31]{31：单侧单位元}\\

1 & \hyperref[ex:1:32]{32：域没有非平凡双边理想}\\

1 & \hyperref[ex:1:33]{33：$Ra$何时是理想}\\

1 & \hyperref[ex:1:34]{34：理想的根}\\

1 & \hyperref[ex:1:35]{35：交换幺环仅有平凡理想时为域}\\

1 & \hyperref[ex:1:36]{36：全矩阵环是单环}\\

1 & \hyperref[ex:1:37]{37：$aba=0$判别零因子}\\

1 & \hyperref[ex:1:38]{38：$1-a$的有限几何级数逆}\\

1 & \hyperref[ex:1:39]{39：交换环的幂零元构成理想}\\

1 & \hyperref[ex:1:40]{40：有限幺环的单侧逆也是双侧逆}\\

1 & \hyperref[ex:1:41]{41：单侧逆不成双侧逆时有无穷多个右逆}\\

1 & \hyperref[ex:1:42]{42：唯一的广义逆推出体}\\

1 & \hyperref[ex:1:43]{43：连续函数环}\\

1 & \hyperref[ex:1:44]{44：有限域上的矩阵计数}\\

1 & \hyperref[ex:1:45]{45：华罗庚恒等式}\\

2 & \hyperref[ex:2:1]{1：阶互素时落入正规子群}\\

2 & \hyperref[ex:2:2]{2：互素阶的分解}\\

2 & \hyperref[ex:2:3]{3：互素指数的开方}\\

2 & \hyperref[ex:2:4]{4：$ab$与$ba$同阶}\\

2 & \hyperref[ex:2:5]{5：给定阶元素的奇偶性}\\

2 & \hyperref[ex:2:6]{6：对称群的中心}\\

2 & \hyperref[ex:2:7]{7：有理数加法群的有限生成子群}\\

2 & \hyperref[ex:2:8]{8：有限阶生成元与交换性推出有限性}\\

2 & \hyperref[ex:2:9]{9：由幂集合刻画循环群}\\

2 & \hyperref[ex:2:10]{10：两组对换生成$S_n$}\\

2 & \hyperref[ex:2:11]{11：两个生成元生成$S_n$}\\

2 & \hyperref[ex:2:12]{12：两个生成元生成$A_n$}\\

2 & \hyperref[ex:2:13]{13：长轮换的中心化子}\\

2 & \hyperref[ex:2:14]{14：$S_5$的共轭类及正规子群}\\

2 & \hyperref[ex:2:15]{15：商群中子群的交}\\

2 & \hyperref[ex:2:16]{16：交子群的陪集}\\

2 & \hyperref[ex:2:17]{17：有限指数子群的交}\\

2 & \hyperref[ex:2:18]{18：有限指数给出正规子群或忠实表示}\\

2 & \hyperref[ex:2:19]{19：最小素因子指数的子群正规}\\

2 & \hyperref[ex:2:20]{20：$p^2$阶群的分类}\\

2 & \hyperref[ex:2:21]{21：非交换6阶群}\\

2 & \hyperref[ex:2:22]{22：15阶群}\\

2 & \hyperref[ex:2:23]{23：10阶群}\\

2 & \hyperref[ex:2:24]{24：不同素数乘积阶的群不单}\\

2 & \hyperref[ex:2:25]{25：$p^2q$阶群有正规西罗子群}\\

2 & \hyperref[ex:2:26]{26：两个非交换的有限例子}\\

2 & \hyperref[ex:2:27]{27：双陪集计数}\\

2 & \hyperref[ex:2:28]{28：西罗2子群只有2阶时的指数2子群}\\

2 & \hyperref[ex:2:29]{29：大子集的乘积}\\

2 & \hyperref[ex:2:30]{30：真子群的共轭不能覆盖有限群}\\

2 & \hyperref[ex:2:31]{31：$S_4$的西罗2子群及西罗3子群}\\

2 & \hyperref[ex:2:32]{32：56阶与148阶群不单}\\

2 & \hyperref[ex:2:33]{33：弗拉蒂尼论证}\\

2 & \hyperref[ex:2:34]{34：正规子群与商群的西罗子群}\\

2 & \hyperref[ex:2:35]{35：子群的西罗子群可由共轭后相交得到}\\

2 & \hyperref[ex:2:36]{36：$\GL_n(\F_p)$的西罗$p$子群}\\

2 & \hyperref[ex:2:37]{37：非交换$p^3$阶群的换位子群}\\

2 & \hyperref[ex:2:38]{38：$p$群中的$p$阶正规子群在中心内}\\

2 & \hyperref[ex:2:39]{39：中心商群循环则原群交换}\\

2 & \hyperref[ex:2:40]{40：有限群嵌入某个交错群}\\

2 & \hyperref[ex:2:41]{41：正规子群与换位子群交平凡时在中心内}\\

2 & \hyperref[ex:2:42]{42：可解群的子群与商群仍可解}\\

2 & \hyperref[ex:2:43]{43：$p^2q$阶群可解}\\

2 & \hyperref[ex:2:44]{44：有限$p$群可解}\\

2 & \hyperref[ex:2:45]{45：两个可解商群给出交核的可解商群}\\

2 & \hyperref[ex:2:46]{46：只有两个内自同构时原群交换}\\

2 & \hyperref[ex:2:47]{47：非平凡有限群至少有两个自同构，除了2阶群}\\

2 & \hyperref[ex:2:48]{48：克莱因四元群的自同构群}\\

2 & \hyperref[ex:2:49]{49：$S_3,S_4$为完全群}\\

2 & \hyperref[ex:2:50]{50：非交换单群的自同构群完全}\\

2 & \hyperref[ex:2:51]{51：不存在换位子群为$S_4$的更大群}\\

2 & \hyperref[ex:2:52]{52：3-轮换为$A_n$中的换位子}\\

2 & \hyperref[ex:2:53]{53：$S_4$的两个生成元及关系}\\

2 & \hyperref[ex:2:54]{54：$A_4$的两个生成元及关系}\\

2 & \hyperref[ex:2:55]{55：四元数群的关系与条件}\\

2 & \hyperref[ex:2:56]{56：60以下的群及60阶单群}\\

2 & \hyperref[ex:2:57]{57：有限指数子群仍有限生成}\\

3 & \hyperref[ex:3:1]{1：由\texorpdfstring{$1-ab$}{1-ab}可逆推出\texorpdfstring{$1-ba$}{1-ba}可逆}\\
3 & \hyperref[ex:3:2]{2：右逆不唯一、零因子与非单位的等价性}\\
3 & \hyperref[ex:3:3]{3：右逆不唯一便有无穷多个右逆}\\
3 & \hyperref[ex:3:4]{4：幂零元全体是否为理想}\\
3 & \hyperref[ex:3:5]{5：素特征的幂公式与一般环的特征}\\
3 & \hyperref[ex:3:6]{6：矩阵环的全部双边理想}\\
3 & \hyperref[ex:3:7]{7：理想积的分配律}\\
3 & \hyperref[ex:3:8]{8：满同态下的元素性质与整环性质}\\
3 & \hyperref[ex:3:9]{9：有理系数四元数构成子体}\\
3 & \hyperref[ex:3:10]{10：实四元数体的中心}\\
3 & \hyperref[ex:3:11]{11：与四元数基元不交换的元素}\\
3 & \hyperref[ex:3:12]{12：共轭不变的四元数子体}\\
3 & \hyperref[ex:3:13]{13：互素理想与理想积的包含}\\
3 & \hyperref[ex:3:14]{14：互素理想之积包含其交中的理想}\\
3 & \hyperref[ex:3:15]{15：只有恒等反自同构的矩阵子环}\\
3 & \hyperref[ex:3:16]{16：整数下三角矩阵环的全部理想}\\
3 & \hyperref[ex:3:17]{17：剩余类矩阵环的单位群阶数}\\
3 & \hyperref[ex:3:18]{18：有限无零因子交换环自动有幺且为域}\\
3 & \hyperref[ex:3:19]{19：理想关系不具传递性}\\
3 & \hyperref[ex:3:20]{20：非单位属于某个极大理想}\\
3 & \hyperref[ex:3:21]{21：极大理想之交的单位刻画}\\
3 & \hyperref[ex:3:22]{22：强拟正则元全体与雅各布森根}\\
3 & \hyperref[ex:3:23]{23：有限交换环的素理想必极大}\\
3 & \hyperref[ex:3:24]{24：整数子群定理与循环群分类的等价性}\\
3 & \hyperref[ex:3:25]{25：单位根群的满而非单自同态}\\
3 & \hyperref[ex:3:26]{26：若尔当同态的幂与极化恒等式}\\
3 & \hyperref[ex:3:27]{27：若尔当同态对一对元素的二择一性质}\\
3 & \hyperref[ex:3:28]{28：逐对二择一推出全局同态或反同态}\\
3 & \hyperref[ex:3:29]{29：除环若尔当同态的乘法方向}\\
3 & \hyperref[ex:3:30]{30：模素数幂环的零因子与素理想}\\
3 & \hyperref[ex:3:31]{31：模素数幂系数多项式的单位与零因子}\\
3 & \hyperref[ex:3:32]{32：直和环的理想分解及整数商环的根与单位群}\\
3 & \hyperref[ex:3:33]{33：三次多项式在六元环中的六个根}\\
3 & \hyperref[ex:3:34]{34：四元数中的二次方程有无穷多个根}\\
3 & \hyperref[ex:3:35]{35：两个平方根之和的极小多项式}\\
3 & \hyperref[ex:3:36]{36：三次代数关系下的乘法和求逆}\\
3 & \hyperref[ex:3:37]{37：八元域的商环构造}\\
3 & \hyperref[ex:3:38]{38：五元域上的全部首一不可约二次式}\\
3 & \hyperref[ex:3:39]{39：二十五元域的构造}\\
3 & \hyperref[ex:3:40]{40：整数多项式环中的一个极大理想}\\
3 & \hyperref[ex:3:41]{41：模三约化诱导的商环同构}\\
3 & \hyperref[ex:3:42]{42：多项式商环无零因子的判据}\\
3 & \hyperref[ex:3:43]{43：幂零多项式的每个系数均幂零}\\
3 & \hyperref[ex:3:44]{44：多项式零因子的常数湮灭元判据}\\
3 & \hyperref[ex:3:45]{45：常数项可逆且其余系数幂零的多项式可逆}\\
3 & \hyperref[ex:3:46]{46：各系数幂零推出多项式幂零}\\
3 & \hyperref[ex:3:47]{47：多元多项式环的幂零根}\\
3 & \hyperref[ex:3:48]{48：有限环中右逆必为左逆}\\
3 & \hyperref[ex:3:49]{49：多项式环中强拟正则元就是幂零元}\\
3 & \hyperref[ex:3:50]{50：每个元素满足某个幂等关系时素理想极大}\\
3 & \hyperref[ex:3:51]{51：布尔环的特征、素商与有限生成理想}\\
3 & \hyperref[ex:3:52]{52：唯一素理想、非单位幂零与根商域的等价性}\\
3 & \hyperref[ex:3:53]{53：只有有限多个理想的整环是域}\\
3 & \hyperref[ex:3:54]{54：理想集合有极小元的整环是域}\\
3 & \hyperref[ex:3:55]{55：华罗庚恒等式中的可逆性与因子次序}\\
3 & \hyperref[ex:3:56]{56：共轭不变除子环的中心性}\\
3 & \hyperref[ex:3:57]{57：实多项式函数环到实数的同态}\\
3 & \hyperref[ex:3:58]{58：无限域上非零多项式的非零点集}\\
3 & \hyperref[ex:3:59]{59：有限域上原点的指示多项式}\\
3 & \hyperref[ex:3:60]{60：只有原点为零点的多项式次数下界}\\
3 & \hyperref[ex:3:61]{61：低于变量数次数的多项式还有非零零点}\\
3 & \hyperref[ex:3:62]{62：形式幂级数的逆、整环性与商域表示}\\
3 & \hyperref[ex:3:63]{63：非单位构成唯一极大理想的条件}\\
3 & \hyperref[ex:3:64]{64：单位加零构成子域及其外元素超越性}\\
3 & \hyperref[ex:3:65]{65：无限域上的乘法型加法恒等式无非常数多项式解}\\
3 & \hyperref[ex:3:66]{66：全矩阵环的左理想都是主左理想}\\
3 & \hyperref[ex:3:67]{67：阿贝尔群的自同态环}\\
3 & \hyperref[ex:3:68]{68：循环群、有理数群与初等阿贝尔群的自同态环}\\
\bottomrule\end{longtable}\end{document}
